\section{AIME II 2025}
        \begin{enumerate}
        	\item (\textit{AIME II 2025}) Six points $A, B, C, D, E,$ and $F$ lie in a straight line in that order. Suppose that $G$ is a point not on the line and that $AC=26$, $BD=22$, $CE=31$, $DF=33$, $AF=73$, $CG=40$, and $DG=30.$ Find the area of $\triangle BGE$.


 
 \textbf{Answer: }468



\textbf{Solution 1}


                    \begin{figure}[H]
                        \centering
                        \includegraphics[width=0.5\textwidth]{Solution/AIME_II_2025_1_Solution_Figure_1.png}
                    \end{figure}
                    

 Let $AB=a$, $BC=b$, $CD=c$, $DE=d$ and $EF=e$. Then we know that $a+b+c+d+e=73$, $a+b=26$, $b+c=22$, $c+d=31$ and $d+e=33$. From this we can easily deduce $c=14$ and $a+e=34$ thus $b+c+d=39$. Using Heron's formula we can calculate the area of $\triangle{CGD}$ to be $\sqrt{(42)(28)(12)(2)}=168$, and since the base of $\triangle{BGE}$ is $\frac{39}{14}$ of that of $\triangle{CGD}$ $\triangle{BGE}$ shares an altitude with $\triangle{CGD}$, we conclude they are proportional and we can calculate the area of $\triangle{BGE}$ to be $168\times \frac{39}{14}=\boxed{468}$.

 





\textbf{Solution 2 (Law of Cosines)}

We need to solve for the lengths of $AB$, $BC$, $CD$, $DE$, and $EF$.  Let $AB = a$, $BC = b$, $CD = c$, $DE = d$, and $EF = e$.  We are given the following system of equations:

 
 \[a + b = 26, \quad b + c = 22, \quad c + d = 31, \quad d + e = 33, \quad a + b + c + d + e = 73.\]

 Substituting $a + b = 26$ and $d + e = 33$ into the equation $a + b + c + d + e = 73$, we get:

 
 \[c = 14.\]

 Thus, we have:

 
 \[a = 18, \quad b = 8, \quad c = 14, \quad d = 17, \quad e = 16.\]

 Next, consider triangle $CDG$, where $CD = 14$, $CG = 40$, and $DG = 30$.  By the Law of Cosines, we have:

 
 \[CD^2 = CG^2 + DG^2 - 2 \times CG \times DG \times \cos(\angle CGD).\]

 Substituting the known values:

 
 \[14^2 = 40^2 + 30^2 - 2 \times 40 \times 30 \times \cos(\angle CGD).\]

 Simplifying:

 
 \[196 = 1600 + 900 - 2400 \cos(\angle CGD).\]

 
 \[2400 \cos(\angle CGD) = 2500 - 196 = 2304.\]

 
 \[\cos(\angle CGD) = \frac{24}{25}.\]

 Therefore, we can find $\sin(\angle CGD)$ using the identity $\sin^2 \theta + \cos^2 \theta = 1$:

 
 \[\sin(\angle CGD) = \sqrt{1 - \left(\frac{24}{25}\right)^2} = \frac{7}{25}.\]

 Now, the area of triangle $CDG$ is

 
 \[\frac{1}{2} \times 40 \times 30 \times \frac{7}{25} = 168.\]

 Noting that the height of triangle $CDG$ is the same as the height of triangle $BGE$, the ratio of the areas of the two triangles will be the same as the ratio of their corresponding lengths. Therefore, the answer is

 
 \[\frac{168 \times 39}{14} = \boxed{468}.\]

 (Feel free to add or correct any LaTeX and formatting)

 





\textbf{Solution 3 (planar cutting, similar to 1, not using ratios however)}


                    \begin{figure}[H]
                        \centering
                        \includegraphics[width=0.5\textwidth]{Solution/AIME_II_2025_1_Solution_Figure_2.png}
                    \end{figure}
                    

 Consider the plane $\mathcal{P}$ that contains line segment $AF$ and point $G$. Note that $AF = 73$. This problem is very lenient, as it directly gives us $AC$ and $CE$, who sum to 57, and therefore $EF = 16$. Also, $BD$ and $DF$ form a segment of length 55, so $AB = 18$. It follows that then because $DF = 33$ and $EF = 16$, $DE = 17$. $CE = 31$, and so $CD = 14$. $BD = 22$, so $BC = 8$.

 
                    \begin{figure}[H]
                        \centering
                        \includegraphics[width=0.5\textwidth]{Solution/AIME_II_2025_1_Solution_Figure_3.png}
                    \end{figure}
                    

 Triangle \( \triangle  GCD \), has a height $h$ that is shared by triangle $\triangle BGE$. The area of \( \triangle  GCD \) through Heron's Formula is 168 square units. One can quickly find the height to be 24. The area of $\triangle BGE$ is then $\frac{1}{2} \cdot 39 \cdot 24 = 39 \cdot 12 = \boxed{468}$.

 




	\item (\textit{AIME II 2025}) Find the sum of all positive integers $n$ such that $n+2$ divides the product $3(n+3)(n^2+9).$


 
 \textbf{Answer: }049



\textbf{Solution 1}

$\frac{3(n+3)(n^{2}+9) }{n+2} \in \mathbb{Z}$

 $\Rightarrow \frac{3(n+2+1)(n^{2}+9) }{n+2} \in \mathbb{Z}$

 $\Rightarrow \frac{3(n+2)(n^{2}+9) +3(n^{2}+9)}{n+2} \in \mathbb{Z}$

 $\Rightarrow 3(n^{2}+9)+\frac{3(n^{2}+9)}{n+2} \in \mathbb{Z}$ 

 $\Rightarrow \frac{3(n^{2}-4+13)}{n+2} \in \mathbb{Z}$

 $\Rightarrow \frac{3(n+2)(n-2)+39}{n+2} \in \mathbb{Z}$

 $\Rightarrow 3(n-2)+\frac{39}{n+2} \in \mathbb{Z}$

 $\Rightarrow \frac{39}{n+2} \in \mathbb{Z}$

 Since $n + 2$ is positive, the positive factors of $39$ are $1$, $3$, $13$, and $39$.

 Therefore, $n = -1$, $1$, $11$ and $37$.

 Since $n$ is positive, $n = 1$, $11$ and $37$.

 $1 + 11 + 37 = \framebox{049}$ is the correct answer

 ～\href{https://artofproblemsolving.com/wiki/index.php/User:Tonyttian}{Tonyttian}

 





\textbf{Solution 2}

We observe that $n+2$ and $n+3$ share no common prime factor, so $n+2$ divides $3(n+3)(n^2+9)$ if and only if $n+2$ divides $3(n^2+9)$.

 By dividing $\frac{3(n^2+9)}{n+2}$ either with long division or synthetic division, one obtains $3n-6+\frac{39}{n+2}$. This quantity is an integer if and only if $\frac{39}{n+2}$ is an integer, so $n+2$ must be a factor of 39. As in Solution 1, $n \in \{1,11,37\}$ and the sum is $\boxed{049}$.

 \begin{verbatim}
NOTE: If you don't observe that it divides if and only if $n+2$ divides $3(n^2+9)$ you can still do synthetic division by expanding $3(n+3)(n^2+9)$ into $3n^3+9n^2+27n+81$ and get the same remainder. -unhappyfarmer
\end{verbatim}





\textbf{Solution 3 (modular arithmetic)}

Let's express the right-hand expression in terms of mod $n + 2$.

 $3 \equiv 3 \pmod{n + 2}$.

 $n + 3 \equiv 1 \pmod{n + 2}$.

 $n^2 + 9 \equiv 13 \pmod{n + 2}$ since $n^2 - 4 \equiv 0 \pmod{n + 2}$ with a quotient $n - 2$

 $3(n + 3)(n^2 + 9) \equiv 3(1)(13) \pmod{n + 2} \equiv 39 \pmod{n + 2}$

 This means $39 = (n + 2)k \pmod{n + 2}$ where k is some integer.

 Note that $n + 2$ is positive, meaning $n + 2 \geq 3$.

 $n + 2$ is one of the factors of 39, so $n + 2 = 3, 13,$ or $39$, so $n = 1, 11,$ or $37$.

 The sum of all possible $n$ is $1 + 11 + 37 = \boxed{049}$.

 





\textbf{Solution 4(Polynomial Remainder Theorem) fast and easy}

For $n+2$ to be a divisor, the remainder needs to be divisible by $n+2$. Using polynomial remainder theorem, we can plug in $-2$ to get $39$ as our remainder. $\frac{39}{n+2}$ is an integer only when $n = 1, 11, 37$. Hence, the sum of all possible $n$ is $1 + 11 + 37 = \boxed{049}$.

 





\textbf{Solution 5 (Synthetic Division on the Entire Polynomial)}


 \[3(n+3)(n^2+9)\]
 \[3(n^3 +3n^2 +9n + 27)\]
 \[3n^3 + 9n^2 + 27n + 81\].
 \[n+2 \ | \ 3n^3 + 9n^2 + 27n + 81\]
 \[\begin{array}{r|rrrr} -2 & 3 & 9 & 27 & 81 \\    &   & -6 & -6 & -42 \\ \hline    & 3 & 3 & 21 & \boxed{39} \end{array}\]

 
 \[n+2 \ | \ 39\]
 \[n + 2 = \not{1}, 3, 13, 39 \ : \ n \in \mathbb{Z}^+\]
 \[n = 1, 11, 37\]
 \[\sum n = 1+11+37 = \boxed{\mathbf{49}}\].

 


 




	\item (\textit{AIME II 2025}) Four unit squares form a $2 \times 2$ grid. Each of the $12$ unit line segments forming the sides of the squares is colored either red or blue in such a way that each unit square has $2$ red sides and $2$ blue sides. One example is shown below (red is solid, blue is dashed). Find the number of such colorings.


 
 \textbf{Answer: }082



\textbf{Solution}

Let the red edges be "used" edges. In the digarams below, dashed lines are uncolored lines yet to be decided. Since all four center edges are common to both squares, we consider five distinct cases:

 $\textbf{Case 1:}$ All center edges are used. There is only one way to do this.

 
                    \begin{figure}[H]
                        \centering
                        \includegraphics[width=0.5\textwidth]{Solution/AIME_II_2025_3_Solution_Figure_1.png}
                    \end{figure}
                    

 $\textbf{Case 2:}$ Three center edges are used, meaning two squares are missing an edge. For each square, there are 2 ways to choose an edge, resulting in $2 \times 2 = 4$ ways. Additionally, considering the rotational symmetry of the arrangement, there are 4 possible rotations, giving a total of $4 \times 4 = 16$ configurations.

 
                    \begin{figure}[H]
                        \centering
                        \includegraphics[width=0.5\textwidth]{Solution/AIME_II_2025_3_Solution_Figure_2.png}
                    \end{figure}
                    

 $\textbf{Case 3:}$ Two center edges are used. There are two sub-cases:

 $\textbf{Scenario 1:}$ The two selected sides are perpendicular to each other. The square diagonally opposite its adjacent square has only one choice, while the other two squares each have two choices. This gives a total of $1 \times 2 \times 2 = 4$ choices. Considering the 4 possible rotations, the total number of configurations is $4 \times 4 = 16$.

 
                    \begin{figure}[H]
                        \centering
                        \includegraphics[width=0.5\textwidth]{Solution/AIME_II_2025_3_Solution_Figure_3.png}
                    \end{figure}
                    

 $\textbf{Scenario 2:}$ The two selected sides are aligned along the same straight line. Each of the four squares has 2 choices, yielding $2^4 = 16$ possible choices. Taking into account the 2 possible rotations, the total number of configurations is $16 \times 2 = 32$.

 
                    \begin{figure}[H]
                        \centering
                        \includegraphics[width=0.5\textwidth]{Solution/AIME_II_2025_3_Solution_Figure_4.png}
                    \end{figure}
                    

 $\textbf{Case 4:}$ Only one center edge is used. This case is similar to Case 2, yielding 16 possible configurations.

 
                    \begin{figure}[H]
                        \centering
                        \includegraphics[width=0.5\textwidth]{Solution/AIME_II_2025_3_Solution_Figure_5.png}
                    \end{figure}
                    

 $\textbf{Case 5:}$ No center edge is used. This is similar to Case 1, with only 1 possible configuration.

 
                    \begin{figure}[H]
                        \centering
                        \includegraphics[width=0.5\textwidth]{Solution/AIME_II_2025_3_Solution_Figure_6.png}
                    \end{figure}
                    

 In conclusion, the total number of configurations is:

 
 \[1 + 16 + 16 + 32 + 16 + 1 = \boxed{\textbf{082}}\]

 




	\item (\textit{AIME II 2025}) The product
 \[\prod^{63}_{k=4} \frac{\log_k (5^{k^2 - 1})}{\log_{k + 1} (5^{k^2 - 4})} = \frac{\log_4 (5^{15})}{\log_5 (5^{12})} \cdot \frac{\log_5 (5^{24})}{\log_6 (5^{21})}\cdot \frac{\log_6 (5^{35})}{\log_7 (5^{32})} \cdots \frac{\log_{63} (5^{3968})}{\log_{64} (5^{3965})}\]is equal to $\frac mn,$ where $m$ and $n$ are relatively prime positive integers. Find $m + n.$


 
 \textbf{Answer: }106



\textbf{Solution 1}

We can rewrite the equation as:

 \begin{align*}&= \dfrac{15}{12} \cdot \dfrac{24}{21} \cdot \dfrac{35}{32} \cdot \dots \cdot \dfrac{3968}{3965} \cdot \dfrac{\log_4{5}}{\log_{64}{5}} \\&= \log_4{64} \cdot \dfrac{(4+1)(4-1)(5+1)(5-1)\cdot \dots \cdot (63+1)(63-1)}{(4+2)(4-2)(5+2)(5-2)\cdot \dots \cdot (63+2)(63-2)} \\&= 3 \cdot \dfrac{5 \cdot 3 \cdot 6 \cdot 4 \cdot \dots \cdot 64 \cdot 62}{6 \cdot 2 \cdot 7 \cdot 3 \cdot \dots \cdot 65 \cdot 61} \\&= 3 \cdot \dfrac{5 \cdot 62}{65 \cdot 2} \\&= 3 \cdot \dfrac{5 \cdot 2 \cdot 31}{5 \cdot 13 \cdot 2} \\&= 3 \cdot \dfrac{31}{13} \\&= \dfrac{93}{13}\end{align*}


 
Desired answer: $93 + 13 = \boxed{106}$

 (Feel free to correct any $\LaTeX$ and formatting.)

 





\textbf{Solution 2}

We can move the exponents to the front of the logarithms like this:\begin{align*}\frac{\log_4 (5^{15})}{\log_5 (5^{12})} \cdot \frac{\log_5 (5^{24})}{\log_6 (5^{21})}\cdot \frac{\log_6 (5^{35})}{\log_7 (5^{32})} \cdots = \frac{15\log_4 (5)}{12\log_5 (5)} \cdot \frac{24\log_5 (5)}{21\log_6 (5)}\cdot \frac{35\log_6 (5)}{32\log_7 (5)} \cdots\end{align*}
Now we multiply the logs and fractions separately. Let's do it for the logs first:\begin{align*}\frac{\log_4 (5)}{\log_5 (5)} \cdot \frac{\log_5 (5)}{\log_6 (5)}\cdot \frac{\log_6 (5)}{\log_7 (5)} \cdots \frac{\log_{63} (5)}{\log_{64} (5)} = \frac{\log_4 (5)}{\log_{64} (5)} = 3\end{align*}
Now fractions:\begin{align*}\frac{15}{12} \cdot \frac{24}{21} \cdot \frac{35}{32} \cdots = \frac{3\cdot 5}{2\cdot 6} \cdot \frac{4\cdot 6}{3\cdot 7} \cdot \frac{5\cdot 7}{4\cdot 8} \cdots \frac{62\cdot 64}{61\cdot 65} = \frac{5}{2} \cdot \frac{62}{65} = \frac{31}{13}\end{align*}
Multiplying these together gets us the original product, which is $\frac{31}{13} \cdot 3 = \frac{93}{13}$.Thus $m+n=\boxed{106}$.

 





\textbf{Solution 3}

Using logarithmic identities and the change of base formula, the product can be rewritten as 
 \[\prod_{k=4}^{63}\frac{k^2-1}{k^2-4}\frac{\log(k+1)}{\log(k)}\]. Then we can separate this into two series. The latter series is a telescoping series, and it can be pretty easily evaluated to be $\frac{\log(64)}{\log(4)}=3$. The former can be factored as $\frac{(k-1)(k+1)}{(k-2)(k+2)}$, and writing out the first terms could tell us that this is a telescoping series as well. Cancelling out the terms would yield $\frac{5}{2}\cdot\frac{62}{65}=\frac{31}{13}$. Multiplying the two will give us $\frac{93}{13}$, which tells us that the answer is $\boxed{106}$.

 





\textbf{Solution 4 (thorough)}

The product is equal to $\prod^{63}_{k=4} \frac{(k-1)(k+1)\log_k 5}{(k-2)(k+2)\log_{k + 1} 5}$ from difference of squares and properties of logarithms. We can now expand:

 \begin{align*}\prod^{63}_{k=4} \frac{(k-1)(k+1)\log_k 5}{(k-2)(k+2)\log_{k + 1} 5} &= \prod^{63}_{k=4}\frac{\log_k 5}{\log_{k+1} 5} \cdot \frac{3 \cdot 5 \cdot 4 \cdot 6 \cdot 5 \cdot 7 \cdots 62 \cdot 64}{2 \cdot 6 \cdot 3 \cdot 7 \cdot 4 \cdot 8 \cdots 61 \cdot 65} \\&= \frac{\log_4 5 \cdot \log_5 5 \cdots \log_{63} 5}{\log_{5} 5 \cdot \log_6 5 \cdots \log_{64} 5} \cdot \frac{3 \cdot 4 \cdot (5^2 \cdot 6^2 \cdots 62^2) \cdot 63 \cdot 64}{2 \cdot 3 \cdot 4 \cdot 5 \cdot (6^2 \cdot 7^2 \cdots 61^2) \cdot 62 \cdot 63 \cdot 64 \cdot 65} \\&= \frac{\log_4 5}{\log_{64} 5} \cdot \frac{3 \cdot 4 \cdot 5^2 \cdot (6^2 \cdots 61^2) \cdot 62^2 \cdot 63 \cdot 64}{2 \cdot 3 \cdot 4 \cdot 5 \cdot (6^2 \cdots 61^2) \cdot 62 \cdot 63 \cdot 64 \cdot 65} \\&= \log_{64} 4 \cdot \frac{3 \cdot 4 \cdot 5 \cdot 62}{2 \cdot 65} \\&= 3 \cdot \frac{31}{13} \\&= \frac{93}{13}\end{align*}


 Thus the answer is $93+13=\boxed{106}$. 

 




	\item (\textit{AIME II 2025}) Suppose $\triangle ABC$ has angles $\angle BAC = 84^\circ, \angle ABC=60^\circ,$ and $\angle ACB = 36^\circ.$ Let $D, E,$ and $F$ be the midpoints of sides $\overline{BC}, \overline{AC},$ and $\overline{AB},$ respectively. The circumcircle of $\triangle DEF$ intersects $\overline{BD}, \overline{AE},$ and $\overline{AF}$ at points $G, H,$ and $J,$ respectively. The points $G, D, E, H, J,$ and $F$ divide the circumcircle of $\triangle DEF$ into six minor arcs, as shown. Find $\widehat{DE}+2\cdot \widehat{HJ} + 3\cdot\widehat{FG},$ where the arcs are measured in degrees.


 
                    \begin{figure}[H]
                        \centering
                        \includegraphics[width=0.5\textwidth]{Problem/AIME_II_2025_5_Problem_Figure_1.png}
                    \end{figure}
                    


 
 \textbf{Answer: }336



\textbf{Solution 1}

Notice that due to midpoints and properties of medial triangle, $\triangle DEF\sim\triangle FBD\sim\triangle AFE\sim\triangle EDC\sim\triangle ABC$ (and by medial triangle $\triangle DEF \cong \triangle FBD \cong \triangle AFE \cong \triangle EDC.$) . As a result, the angles and arcs are readily available. Due to inscribed angles,
 \[\widehat{DE}=2\angle DFE=2\angle ACB=2\cdot36=72^\circ\]Similarly,
 \[\widehat{FG}=2\angle FDB=2\angle ACB=2\cdot36=72^\circ\]

 In order to calculate $\widehat{HJ}$, we use the fact that $\angle BAC=\frac{1}{2}(\widehat{FDE}-\widehat{HJ})$. We know that $\angle BAC=84^\circ$, and
 \[\widehat{FDE}=360-\widehat{FE}=360-2\angle FDE=360-2\angle CAB=360-2\cdot84=192^\circ\]

 Substituting,

 \begin{align*}84 &= \frac{1}{2}(192-\widehat{HJ}) \\168 &= 192-\widehat{HJ} \\\widehat{HJ} &= 24^\circ\end{align*}


 Thus, $\widehat{DE}+2\cdot\widehat{HJ}+3\cdot\widehat{FG}=72+48+216=\boxed{336}^\circ$. 

 Alternatively, 

 \begin{align*} \widehat{HJ} &= \widehat{FH} + \widehat{JE} - \widehat{FE} \\ &= 2\angle FEH + 2\angle JFE - 2\angle FDE \\ &= 2 \cdot 36^\circ + 2 \cdot 60^\circ - 2 \cdot 84^\circ \\ &= 24^\circ. \end{align*}


 





\textbf{Solution 2}

Notice that $\triangle ABC \sim \triangle FHE \sim \triangle DEF$ because $FE \parallel BC$ and $DE \parallel AB,$ so all angles in each triangle will be equal by the Midline Theorem. Therefore, we have $\widehat{DE} = 2 \cdot 36^\circ = 72^\circ.$ Now, quadrilateral $FHED$ is cyclic, so opposite angles add to $180^\circ.$ Since we know from similar triangles that $\angle{FED} = 60^\circ + 36^\circ = 96^\circ,$ we see that $\angle{HFD} = 84^\circ.$ We also know that $\angle{JFE} = 60^\circ + 36^\circ = 96^\circ,$ so that means $\angle{JFH} = 96^\circ - 84^\circ = 12^\circ \implies \widehat{JH} = 24^\circ.$ Finally, $\angle{B} = 60^\circ = \frac{\widehat{JED} - \widehat{FG}}{2}.$ 
 \[\widehat{JED} = \widehat{DE} + \widehat{JE} = 72^\circ + 2(\angle{JFE}) = 192^\circ.\] So $\widehat{FG} = 192^\circ - 120^\circ = 72^\circ.$ The answer is $72^\circ + 2\cdot 24^\circ + 3\cdot 72^\circ = \boxed{336^\circ}.$

 
NOTE: Several of the letters appear to be misused. For example, $\angle{FED}=60^\circ$ not $96^\circ$.

 NOTE2: It should be $\angle{HED}=60^\circ + 36^\circ = 96^\circ$ and $\angle{HFD} = 180^\circ - \angle{HED} = 84^\circ$. The same for $\angle{JFD}$, not $\angle{JFE}$.

 





\textbf{Solution 3}

Triangle $DEF$ is the median triangle of $ABC$. It's known that the circumcircle of $DEF$ is the Nine-Point Circle of $ABC$. 

 Arc $DE$ is easy to find through inscribed angles theorem, as $72^\circ$. 

 Through Intercepted Secants Theorem, the measure of angle $JAH$ is precisely half that of the difference of the far arc $FE$ and the near arc $JH$. Thus, $84 = 0.5 \cdot (192-JH)$, or $24 = JH$

 Finally, through Alternate Interior it's obvious that angle $FDG$ is $36^\circ$. It follows that arc $FG$ is just $72^\circ$.

 
 \[DE + 2HJ + 3FG:~\text{arcs measured in degrees}~ = 72 + 48 + 216 = \boxed{336}^\circ\]

 The purpose of this solution is to bring up the existence of the nine-point circle. Consequently, points $G$, $H$, and $J$ are not random, but are the intersections of the altitudes dropped from $A$, $B$, and $C$ respectively. It may not be used in this problem specifically, but is great information for the long run.

 




	\item (\textit{AIME II 2025}) Circle $\omega_1$ with radius $6$ centered at point $A$ is internally tangent at point $B$ to circle $\omega_2$ with radius $15$. Points $C$ and $D$ lie on $\omega_2$ such that $\overline{BC}$ is a diameter of $\omega_2$ and ${\overline{BC} \perp \overline{AD}}$. The rectangle $EFGH$ is inscribed in $\omega_1$ such that $\overline{EF} \perp \overline{BC}$, $C$ is closer to $\overline{GH}$ than to $\overline{EF}$, and $D$ is closer to $\overline{FG}$ than to $\overline{EH}$, as shown. Triangles $\triangle {DGF}$ and $\triangle {CHG}$ have equal areas. The area of rectangle $EFGH$ is $\frac{m}{n}$, where $m$ and $n$ are relatively prime positive integers. Find $m+n$.


 
                    \begin{figure}[H]
                        \centering
                        \includegraphics[width=0.5\textwidth]{Problem/AIME_II_2025_6_Problem_Figure_1.png}
                    \end{figure}
                    


 
 \textbf{Answer: }293



\textbf{Solution 1 (Thorough)}

Let $GH=2x$ and $GF=2y$. Notice that since $\overline{BC}$ is perpendicular to $\overline{GH}$ (can be proven using basic angle chasing) and $\overline{BC}$ is an extension of a diameter of $\omega_1$, then $\overline{CB}$ is the perpendicular bisector of $\overline{GH}$. Similarly, since $\overline{AD}$ is perpendicular to $\overline{GF}$ (also provable using basic angle chasing) and $\overline{AD}$ is part of a diameter of $\omega_1$, then $\overline{AD}$ is the perpendicular bisector of $\overline{GF}$.

 From the \href{https://artofproblemsolving.com/wiki/index.php?title=Pythagorean\_Theorem}{Pythagorean Theorem} on $\triangle GFH$, we have $(2x)^2+(2y)^2=12^2$, so $x^2+y^2=36$. To find our second equation for our system, we utilize the triangles given.

 Let $I=\overline{GH}\cap\overline{CB}$. Then we know that $GFBI$ is also a rectangle since all of its angles can be shown to be right using basic angle chasing, so $FG=IB$. We also know that $CI+IB=2\cdot 15=30$. $IA=y$ and $AB=6$, so $CI=30-y-6=24-y$. Notice that $CI$ is a height of $\triangle CHG$, so its area is $\frac{1}{2}(2x)(24-y)=x(24-y)$.

 Next, extend $\overline{DA}$ past $A$ to intersect $\omega_2$ again at $D'$. Since $\overline{BC}$ is given to be a diameter of $\omega_2$ and $\overline{BC}\perp\overline{AD}$, then $\overline{BC}$ is the perpendicular bisector of $\overline{DD'}$; thus $DA=D'A$. By Power of a Point, we know that $CA\cdot AB=DA\cdot AD'$. $CA=30-6=24$ and $AB=6$, so $DA\cdot AD'=(DA)^2=24\cdot6=144$ and $DA=D'A=12$.

 Denote $J=\overline{DA}\cap\overline{GF}$. We know that $DJ=DA-AJ=12-x$ (recall that $GI=IH=x$, and it can be shown that $GIAJ$ is a rectangle). $\overline{DJ}$ is the height of $\triangle DGF$, so its area is $\frac{1}{2}(2y)(12-x)=y(12-x)$.

 We are given that $[DGF]=[CHG]$ ($[ABC]$ denotes the area of figure $ABC$). As a result, $x(24-y)=y(12-x)$. This can be simplified to $y=2x$. Substituting this into the Pythagorean equation yields $5x^2=36$ and $x=\frac{6}{\sqrt{5}}$. Then $y=\frac{12}{\sqrt{5}}$.

 $[EFGH]=2x\cdot2y=2\cdot\frac{6}{\sqrt{5}}\cdot2\cdot\frac{12}{\sqrt{5}}=\frac{288}{5}$, so the answer is $288+5=\boxed{293}$. 

 





\textbf{Solution 2 (Faster)}

Denote the intersection of $BC$ and $w_1$ as $P$, the intersection of $BC$ and $GH$ be $Q$, and the center of $w_2$ to be $O$. Additionally, let $EF = GH = a, FG = EH = b$. We have that $CP = 18$ and $PQ = 6-\frac{b}{2}$. Considering right triangle $OAD$, $AD = 12$. Letting $R$ be the intersection of $AD$ and $FG$, $DR = 12 - \frac{a}{2}$. Using the equivalent area ratios: 
 \[\frac{a(24-\frac{b}{2})}{2} = \frac{(12-\frac{a}{2})b}{2}\]

 This equation gives $b=2a$. Using the \href{https://artofproblemsolving.com/wiki/index.php?title=Pythagorean\_Theorem}{Pythagorean Theorem} on triangle $GHE$ gives that $a^2+b^2 = 144$. Plugging the reuslt $b=2a$ into this equation gives that the area of the triangle is $\frac{288}{5} \to \boxed{293}$.

 





\textbf{Solution 3 (Not Recommended)}

You can use your ruler to check that $GF=2GH.$ Then, we have $GF^2+GH^2=36,$ and we solve the system of equations to get $\frac{288}{5},$ so the answer is $\boxed{293}.$

 Note: This method is not recommended as the diagrams are not necessarily drawn to scale. However, it can be used in emergency situations or to verify the answer.

 


 





\textbf{Solution 4 (Almost no Algebra)}

If we draw segments connecting $A$ to $F, G,$ and $H,$ we can easily verify that all of the right triangles created are congruent. Thus, triangles $AGF$ and $AGH$ have equal areas, which means, by the given conditions, that kites $AGCH$ and $AGDF$ have equal areas. Thus, by the area formula for kites, 

 
 \[\frac{1}{2}(AD)(GF)=\frac{1}{2}(AC)(GH),\] or 
 \[(AD)(GF)=(24)(GH).\] 

 Also, if we extend $AD$ to the other side of the large circle, the chord length formula gives 

 
 \[(AD)(AD)=(AC)(AB),\] so 
 \[(AD)(AD)=(24)(6).\] 

 
Squaring the second equation above gives 
 \[(AD)(AD)(GF)(GF)=(24)(24)(GH)(GH).\] Dividing by the fourth equation gives 
 \[(GF)^2=4(GH)^2.\]Since we know, by the Pythagorean Theorem, that $(GF)^2+(GH)^2=144,$ we can substitute to determine that $5(GH)^2=144.$ 

 The desired area is $(GF)(GH)=2(GH)^2=\frac{2}{5}(144),$ which is $\frac{288}{5},$ so the answer is $\boxed{293}.$

 




	\item (\textit{AIME II 2025}) Let $A$ be the set of positive integer divisors of $2025$. Let $B$ be a randomly selected subset of $A$. The probability that $B$ is a nonempty set with the property that the least common multiple of its elements is $2025$ is $\frac{m}{n}$, where $m$ and $n$ are relatively prime positive integers. Find $m+n$.


 
 \textbf{Answer: }237



\textbf{Solution 1}

We split into different conditions:

 Note that the numbers in the set need to have a least common multiple of $2025$, so we need to ensure that the set has at least 1 number that is a multiple of $3^4$ and a number that is a multiple of $5^2$. 

 Multiples of $3^4$: $81, 405, 2025$

 Multiples of $5^2$: $25, 75, 225, 675, 2025$

 If the set $B$ contains $2025$, then all of the rest $14$ factors is no longer important. The valid cases are $2^{14}$, which can be thought of as selecting either "Yes" or "No" ($2$ possibilities) for each of the remaining $14$ factors.

 If the set $B$ doesn't contain $2025$, but contains $405$, we just need another multiple of $5^2$. It could be 1 of them, 2 of them, 3 of them, or 4 of them, which has $2^4 - 1 = 15$ cases. Excluding $2025, 405, 25, 75, 225, 675,$ the rest 9 numbers could appear or not appear. Therefore, this case has a valid case of $15 \cdot 2^9$. 

 If set $B$ doesn't contain $2025$ nor $405$, it must contain $81$. It also needs to contain at least 1 of the multiples from $5^2$, where it would be $15 \cdot 2^8$. 

 The total valid cases are $2^{14} + 15 \cdot (2^9 + 2^8)$, and the total cases are $2^{15}$. 

 The answer is $\cfrac{2^8 \cdot (64 + 30 + 15)}{2^8 \cdot 2^7}= \frac{109}{128}$. 

 Desired answer: $109 + 128 = \boxed{237}$.

 





\textbf{Solution 2}

We take the complement and use PIE. Suppose the LCM of the elements of the set is NOT $2025$. Since $2025=3^4 \cdot 5^2$, it must be that no element $x$ in the subset satisfies $v_3(x)=4$ OR no element $x$ in the subset satisfies $v_5(x)=2$ (in this case, $v_p(n)$ gives the exponent of $p$ in the prime factorization of $n$). For the first case, there are $4 \cdot 3 = 12$ possible divisors that could be in our subset ($v_3(x)=0,1,2,3,v_5(x)=0,1,2$ are possible), for a total count of $2^{12}$ subsets. In the second case, there are $5 \cdot 2 = 10$ possible divisors that could be in our subset, for a total count of $2^{10}$ subsets. However, if both conditions are violated, then there are $4 \cdot 2 = 8$ divisors that could be in our subset, for a total count of $2^8$ subsets. Hence, by PIE, the number of subsets whose LCM is NOT $2025$ is equal to $2^{12}+2^{10}-2^8$. The answer is then 
 \[1-\frac{2^{12}+2^{10}-2^8}{2^{5 \cdot 3}}=\frac{109}{128} \implies \boxed{237}.\]

 





\textbf{Solution 3}

Write numbers in the form of $3^{a}5^{b}$ where $0\leq a\leq 4; 0\leq b\leq 2$

 There are $(4+1)(2+1)=15$ possible divisors of $2025$, so the cardinality of the subsets is $2^{15}$

 If I select $3^4\cdot 5^2$, then I guarantee the LCM is $2025$, so the other $14$ numbers yield $2^{14}$ cases.

 If I select $3^4\cdot 5$, then I must select at least one of $3^a\cdot5^2$, but I can select any other $9$ numbers, so there are 
 \[2^9 \left(\binom{4}{1}+\binom{4}{2}+\binom{4}{3}+\binom{4}{4} \right)=2^9\cdot 15\] ways.

 If I select $3^4$, since we can't select $3^4\cdot 5$ or $3^4\cdot5^2$ anymore, there are $2^8 \left(\binom{4}{1}+\binom{4}{2}+\binom{4}{3}+\binom{4}{4} \right)=2^8\cdot 15$ ways.

 The answer is then $\frac{2^8(15+30+64)}{2^{15}}=\frac{109}{128}\implies \boxed{237}$.

 





\textbf{Solution 4}

We take $2$ cases, depending on whether $2025=3^4\cdot5^2$ is an element of $B$ or not. Here is a table that would be helpful to draw:

 
$\begin{array}{|c|c|c|c|} \hline Divisors  & 5^0=1  & 5^1=5  & 5^2=25  \\ \hline 3^0=1  & 1  & 5  & 25  \\ \hline 3^1=3  & 3 & 15 & 75 \\ \hline 3^2=9 & 9 & 45 & 225 \\ \hline 3^3=27 & 27 & 135 & 675 \\ \hline 3^4=81 & 81 & 405 & 2025 \\ \hline \end{array}$

 
$\textbf{Case 1}$: $2025$ is in $B$.

 If $2025$ is in $B$, then the choices of the rest of the elements doesn't matter, as the least common multiple of the elements of $B$ is already forced to be $2025$. We can either include or not include each of the remaining $14$ possible elements, giving $2$ possibilities for each of the $14$ elements or $2^{14}$ scenarios.

 
$\textbf{Case 2}$: $2025$ is not in $B$.

 If $2025$ is not in $B$, we still must create $B$ such that it includes a term divisible by $5^2$ and a term divisible by $3^4$. Observing the table above, we see that we must choose at least one of $81$ and $405$ (both of which are divisible by $3^4$) and at least one of $25$, $75$, $225$, and $675$ (all of which are divisible by $5^2$). The number of ways to choose at least one of $81$ and $405$ is equal to (ways to choose $1$ of them)$+$(ways to choose $2$ of them)$=\binom{2}{1}+\binom{2}{2}=2+1=3$. The number of ways to choose at least one of $25$, $75$, $225$, and $675$ is equal to (ways to choose $1$ of them)$+$(ways to choose $2$ of them)$+$(ways to choose $3$ of them)$+$(ways to choose $4$ of them)$=\binom{4}{1}+\binom{4}{2}+\binom{4}{3}+\binom{4}{4}=4+6+4+1=15$. We can now choose any number of the remaining $15-1-2-4=8$ divisors of $2025$ (since there are already at least $2$ elements in $B$, we can choose $0$ more divisors without creating the $\varnothing$; since $2025$ is not included in this case, we can choose all $8$ divisors without repeating the scenario from Case $1$ where all the elements of $A$ are in $B$). The number of ways to do so is (ways to choose $1$ of them)$+$(ways to choose $2$ of them)$+$(ways to choose $3$ of them)$+...+$(ways to choose $8$ of them)$=\binom{8}{0}+\binom{8}{1}+\binom{8}{2}+...+\binom{8}{8}=2^8$. Multiplying these independent ways yields $3\cdot15\cdot2^8$ scenarios.

 


 The total number of possible subsets of $A$ (including the $\varnothing$ and $A$ itself) is equal to $\binom{15}{0}+\binom{15}{1}+\binom{15}{2}+...+\binom{15}{15}=2^{15}$.

 


 Therefore, the desired probability is $\frac{2^{14}+3\cdot15\cdot2^8}{2^{15}}=\frac{2^{6}+3\cdot15}{2^{7}}=\frac{64+45}{128}=\frac{109}{128}$, so $m+n=109+128=\boxed{237}$.

 




	\item (\textit{AIME II 2025}) From an unlimited supply of $1$-cent coins, $10$-cent coins, and $25$-cent coins, Silas wants to find a collection of coins that has a total value of $N$ cents, where $N$ is a positive integer. He uses the so-called $\textit{greedy algorithm}$, successively choosing the coin of greatest value that does not cause the value of his collection to exceed $N$. For example, to get $42$ cents, Silas will choose a $25$-cent coin, then a $10$-cent coin, then $7$ $1$-cent coins. However, this collection of $9$ coins uses more coins than necessary to get a total of $42$ cents; indeed, choosing $4$ $10$-cent coins and $2$ $1$-cent coins achieves the same total value with only $6$ coins. 


 In general, the greedy algorithm succeeds for a given $N$ if no other collection of $1$-cent, $10$-cent, and $25$-cent coins gives a total value of $N$ cents using strictly fewer coins than the collection given by the greedy algorithm. Find the number of values of $N$ between $1$ and $1000$ inclusive for which the greedy algorithm succeeds.


 
 \textbf{Answer: }610



\textbf{Solution 1}

We begin by noting that all values of $N \leq 25$ work without issue.

 Starting from $N = 25$ to $29$, the greedy algorithm will select the $25$-cent coin, and no problem arises.

 From $N = 30$ to $34$, the greedy algorithm will select the $25$-cent coin along with $5$ $1$-cent coins to reach a total of $30$, while the optimal solution would involve using $3$ $10$-cent coins. This issue is resolved from $N = 35$ to $39$, as the greedy algorithm can now select $25 + 10$-cent coins to match the optimal solution.

 From $N = 40$ to $44$, a similar problem occurs again. The greedy algorithm selects $25 + 10 + 5 \times 1$-cent coins to reach 40, while the optimal solution would use 4 $10$-cent coins.

 The problem occurs again from $N = 55$ to $59$, where $50 + 5 \times 1$ is not as good as using $25 + 3 \times 10$, and it is resolved at $N = 60$. From $N = 65$ to $69$, a similar issue arises, as $25 \times 2 + 10 + 5 \times 1$ is not as optimal as $25 + 4 \times 10$ to approach 65.

 We observe that this issue repeats in cycles of $25$ numbers, with $10$ of the $25$ numbers in each cycle not working. The cycle starts at $30$, and the next cycle will start $25$ numbers later, at $55$, then $80$, and so on, continuing until $980$–$1005$ for the last cycle.

 The total number of cycles is given by:

 
 \[\frac{955 - 30}{25} + 1 = 38,\]

 and each cycle contains $10$ problematic numbers. Therefore, the total number of problematic numbers is:

 
 \[38 \times 10 = 380.\]

 The cycle from $980$ to $1005$ has the problematic numbers from $980$ to $984$ and $990$ to $994$, giving another $10$ problematic numbers.

 Thus, the total number of unsuccessful numbers from $1$ to $1000$ inclusive is $390$, and the desired count of successful numbers is:

 
 \[1000 - 390 = \boxed{610}.\]

 (Feel free to make any edits on grammar, $\LaTeX$ and formatting.)

 \textbf{Proof of the Existence of the Cycle} For completeness, we'll present a proof that the cycle described in the first solution exists. 

 For a positive integer $n$ in the form $n=30+25m$ where $m$ is a non-negative integer, we claim that in the list of $25$ numbers below, 
 \[\underbrace{n,\ldots,n+4}_{\text{succeeds}},\underbrace{n+5,\ldots, n+9}_{\text{does not succeed}}, \underbrace{n+10,\ldots,n+14}_{\text{does not succeed}}, \underbrace{n+15, \ldots, n+24}_{\text{succeeds}}\] is the situation. We proceed by induction to show that this is true. 

 The base case of $n=30$ or equivalently $m=0$ is shown in Solution 1. Then, if the claim is true for $m=k$ for some non-negative integer $k$, then consider that for $m=k+1$ (i.e $n$ is increased by $25$), the greedy algorithm will first choose a $25$ cent coin for each of the $25$ numbers in the list, which reduces the problem to finding the numbers that succeeds for $m=k$. Hence, by induction, the claim is true.

 \href{https://mathproblemsolvingskills.wordpress.com/}{MathProblemSolvingSkills.com}

 


 




	\item (\textit{AIME II 2025}) There are $n$ values of $x$ in the interval $0<x<2\pi$ where $f(x)=\sin(7\pi\cdot\sin(5x))=0$. For $t$ of these $n$ values of $x$, the graph of $y=f(x)$ is tangent to the $x$-axis. Find $n+t$.


 
 \textbf{Answer: }149



\textbf{Solution 1}

For $\sin(7\pi\cdot\sin(5x))=0$ to happen, whatever is inside the function must be of form $k\pi$. We then equate to have\begin{align*}7\pi\cdot\sin(5x)=k\pi \\\sin(5x)=\frac{k}{7} \\\end{align*}
We know that $-1\le \sin{5x} \le 1$, so clearly $k$ takes all values $-7\le k \le 7$. Since the graph of $\sin{5x}$ has 5 periods between $0$ and $360$, each of the values $k=-6,-5,-4...-1,1,2...6$ give $10$ solutions each. $k=-7,7$ give $5$ solutions each and $k=0$ gives $9$ solutions (to verify this sketch a graph). Thus, $n=139$.

 We know that the function is tangent to the x-axis if it retains the same sign on both sides of the function. This is not true for points at $k=-6,-5,-4...4,5,6$ because one side will be positive and one will be negative. However this will happen if $k=-7,7$ because the sine function "bounces back" and goes over the same values again (to clarify, because $y=\sin{5x}$ has the same values on both sides of $\sin{5x}=\frac{k}{7}$ for $k=-7,7$, as it's a maximum $\sin{5x}=\frac{7}{7}=1$ or minimum $\sin{5x}=\frac{-7}{7}=-1$, the corresponding values of $y=\sin(7\pi\cdot\sin(5x))$ also have a maximum/minimum and repeat around those $x$-values), and $t=10$ of these values exist.Thus, $n+t=\boxed{149}$.

 \textbf{Note} The graph of $7\sin(5x)$ has minimums and maximums when $k = -7,7$ and otherwise crosses so when we take the sine of that interval all the minimums and maximums of the inscribed function "fold" to the $x$ axis.

 





\textbf{Solution 2 (Calculus)}

For $f(x)=0$, we must have $7\pi\cdot\sin(5x)=k\pi$ for some integer $k$. Then $\sin(5x)=\frac{k}{7}$ always satisfies the equation. Notice that on each period of $\sin(5x)$, each $k\in\{-6,-5,\ldots,5,6\}$ is a $y$-value at two distinct points, and each $k=\pm7$ is a $y$-value at one point each. Thus each period has $13\cdot2+2\cdot1=28$ points satisfying the equation. Since the period is $\frac{2\pi}{5}$ and the domain has a length of $2\pi$, we find that $5$ periods occur in our domain if we include $x=0,2\pi$. Adding the case where $x=0$, there are a total of $28\cdot5+1=141$ roots over $x\in[0,2\pi]$. Subtracting the cases at $x=0$ and $x=2\pi$ yields $139$ total roots. This is our $n$.

 
Next, we take the derivative of $f(x)$; using a hideous combination of chain rules we find that

 
 \[f'(x)=35\pi\cos(5x)\cos(7\pi\sin(5x))=0\]

 Thus, for a point to be tangent to the $x$-axis, we must have either $\cos(5x)=0$ or $\cos(7\pi\sin(5x))=0$. In the first case, we know that $\sin(5x)=\frac{k}{7}$ from earlier, so $\cos(5x)=\sqrt{1-\left(\frac{k}{7}\right)^2}=0$. Then $\left(\frac{k}{7}\right)^2=1$, so $k=\pm7$. Recall that over each of the five periods, only one point will satisfy $k=7$, and only one point will satisfy $k=-7$. Thus there are $2\cdot5=10$ points in this case.

 In the second case, we must have $\cos(7\pi\sin(5x))=0$. Substituting $\sin(5x)=\frac{k}{7}$ yields $\cos(k\pi)=0$. But this is impossible since $\cos(0)=1$ and $\cos(\pi)=-1$, so there are no points in this case.

 As a result, $t=10+0=10$, so $n+t=139+10=\boxed{149}$. 

 




	\item (\textit{AIME II 2025}) Sixteen chairs are arranged in a row. Eight people each select a chair in which to sit so that no person sits next to two other people. Let $N$ be the number of subsets of $16$ chairs that could be selected. Find the remainder when $N$ is divided by $1000$.


 
 \textbf{Answer: }907



\textbf{Solution 1 (Recursion)}

Notice that we can treat each chair as an empty space. If a person selects a chair, we fill in the corresponding space with a '$1$': otherwise, we fill in the corresponding space with a '$0$'. Since no person can sit to two other people (and two other people means having a person to your left and your right), that means that we can't have three people sitting in a row). So now the problem boils down to the following: we have a string of $0$'s and $1$'s of length $16$, $8$ of which will be $1$ and $8$ of which will be $0$. This string cannot have three consecutive $1$'s in a row. We need to find the number of ways to construct this string.

 Let's try recursion. The motivation for recursion is that $16$ and $8$ are relatively big numbers (the upper bound is $16$ choose $8$ which is over $12000$!) Also, many problems involving strings and counting with restrictions like 'three in a row' can usually be solved by breaking it up into smaller problems. 

 Define $S(n,k)$ to be the number of ways to create a string of length $n$ with $k$ $1$'s and $n-k$ $0$'s such that we can't have three consecutive $1$'s. We need to find $S(16,8)$. 

 So now we're going to take our recursive step, which usually involves splitting the string into two smaller strings. Suppose we split the string into two strings of length $n-b$ and $b$. Also suppose that in the string of length $b$ there are $i$ $1$'s (which means that in the string of length $n-b$, there are $k-i$ $1$'s. Now consider the first two digits in the string of length $b$: we do casework on them.

 $\textbf{Case 1: 00}$

 In this case, we can imagine the string as \_ \_ \_ \_ ... \_ \_ | 0 0 \_ \_ ... \_ \_ where the vertical line | shows the split. To the left of the split, we can choose the string without constraints in $S(n-b, k-i)$ ways since there are $n-b$ spaces and $k-i$ $1$'s to place there. To the right of the $0$, there are $b-2$ spaces and we still have to place $i$ $1's$, so this can be done in $S(b-2,i)$ ways. Now choosing the left string and the right string are independent, so the number of ways in this case is 
 \[S(n-b, k-i)S(b-2, i)\].

 $\textbf{Case 2: 01}$

 Now the split string looks something like \_ \_ \_ \_ ... \_ \_ | 0 1 \_ \_ ... \_ \_ . The number of ways to choose the string to the left remains unchanged: $S(n-b, k-i)$ ways. Similarly, it would seem that the number of ways to choose the remainder of the string on the right is $S(b-2, i-1)$ (just as in the previous case but we've used up one of the $i$ $1$'s already). However we're overcounting, since if the remainder of the string on the right starts with 1 1 0 \_ \_ ... it makes a valid string by itself, but together with the 0 1 at the start, it makes 0 1 1 1 0, which is not allowed. Therefore we subtract the number of ways that we start with 0 1 1 1 0: there are $b-5$ spaces left and $i-3$ $1$'s left to allocate, so we subtract $S(b-5, i-3)$. So the total in this case is 
 \[S(n-b,k-i) \cdot (S(b-2, i-1)-S(b-5, i-3))\]

 $\textbf{Case 3: 10}$

 Now the split string looks something like \_ \_ \_ \_ ... \_ \_ | 1 0 \_ \_ ... \_ \_ . Now the number of ways to fill in the empty slots to the right is unconstrained as it's following a $0$: there are $b-2$ spaces and $i-1$ $1$'s to allocate, so this gives $S(b-2, i-1)$ ways. Now for the left, we overcount just like in the previous case: filling $k-i$ $1$'s in $n-b$ slots can be done in $S(n-b, k-i)$ ways. However any way ending with 0 1 1 violates the three $1$'s in a row rule since in the center we have 0 1 1 1 0. The number of violations is the number of ways to fill in the empty slots to the left of \_ \_ \_ ... \_ \_ 0 1 1 |, which can be done in $S(n-b-3, k-i-2)$ ways since we've already allocated $2$ $1$'s and used up $3$ spaces. Therefore the total in this case is:

 
 \[S(b-2, i-1) \cdot (S(n-b, k-i)-S(n-b-3, k-i-2))\]

 $\textbf{Case 4: 11}$

 In this case, we have something like \_ \_ \_ \_ ... \_ \_ | 1 1 \_ \_ ... \_ \_ . Since we can't have three $1$'s in a row, we can automatically fill $0$'s in: \_ \_ \_ \_ ... \_ 0 | 1 1 0 \_ ... \_ \_ . Now the number of ways to fill in the spaces on the left is $S(n-b-1, k-i)$ and the number of ways to fill in the spaces on the right is $S(b-3, i-2)$. So the number of ways for this case is

 
 \[S(n-b-1,k-i)S(b-3,i-2)\]

 


 Now that we've finished casework, note that $i$ can be anything: at minimum $i$ can be zero, meaning that the string on the right has $0$ $1$'s, and at maximum $i$ can be $b$: the string is filled with $1$'s. To encapsulate all of the cases of $i$, we sum over the values of $i$:

 
 \[S(n,k) = \sum_{i = 0}^{b} S(n-b, k-i)S(b-2, i) + S(n-b,k-i) \cdot (S(b-2, i-1)-S(b-5, i-3))\]
 \[+ S(b-2, i-1) \cdot (S(n-b, k-i)-S(n-b-3, k-i-2)) + S(n-b-1,k-i)S(b-3,i-2)\]

 
This is our recurrence relationship. But where do we make the split? The above formula works for any $b$, but if we make the split too close to the end, we could end up having to make lots of computations: if we split the string into $15$ and $1$, we'll have to compute $S(15,8)$, $S(14,8)$, etc. which could end up being a lot of work. We want to choose the best $b$ to minimize the number of computations to save time. If we split the string unequally, one side will be larger and that side will require more computations. Therefore, we split the string in half so we just need to compute values of $n$ less than $8$.

 So suppose $b = 8$ and plug in $n=16$, $k=8$. Then the above recurrence relation becomes 

 
 \[S(16,8) = \sum_{i = 0}^{8} S(8, 8-i)S(6, i) + S(8,8-i) \cdot (S(6, i-1)-S(3, i-3))\]
 \[+ S(6, i-1) \cdot (S(8, 8-i)-S(5, 6-i)) + S(7,8-i)S(5,i-2)\]

 
Before we start with computations, let's state our base cases for recursion. Note that if $n<0$ and $k = 0$, $S(n,k) = 1$. Otherwise if $n<0$ or $k<0$ or $k>n$, $S(n,k) = 0$. Also, $S(x,0) = 1$, $S(x,1) = x$, $S(x,2) = \binom{x}{2}$, $S(x,3) = \binom{x}{3}-x+2$ (the last one subtracting $x-2$ for the number of three $1$'s in a row possible). Finally, for other cases of $S(n,k)$ we don't want to use the earlier large recursive formula since we'll have to break those into many small components. Instead, we modify the formula so that $b=2$. This should be fine as $n$ has already been broken down to at most $8$, so these computations shouldn't balloon. Using $b=2$, we get $S(n,k) = S(n-2,k) + 2S(n-2,k-1) - S(n-5,k-3)+S(n-3,k-2)$.

 
Let's start the computation through the earlier cases. 

 $\textbf{Case 1}$

 
 \[\sum_{i = 0}^{8} S(8,8-i)S(6,i) = \sum_{i = 2}^{4} S(8,8-i)S(6,i)\]

 This is because for $i = 5$, $S(6,5) = 0$ since the best we can do without three in a row is 1 1 0 1 1 0, which is four $1$'s, and it's impossible to achieve $5$ $1$'s without three back to back. Also if $i=5$ makes $S(6,i)$ $0$, then $i \geq 5$ makes $S(6,i)$ $0$. Similarly we can't fit $7$ $1$'s in $8$ slots without three back-to-back, so $i$ can't be $0$ or $1$. So $i$ must go from $2$ to $4$. In future cases, I'll skip the explanation of why we tighten our bounds on $i$ since the reasoning is the same as here.

 
 \[\implies \sum_{i = 0}^{8} S(8,8-i)S(6,i) = S(8,6)S(6,2) + S(8,5)S(6,3) + S(8,4)S(6,4)\]

 Applying the modified recurrence relation for $b=2$, $S(8,6) = S(6,6) + 2S(6,5)-S(3,3)+S(5,4) = 0+0-0+S(5,4)$$= S(3,4)+2S(3,3)-S(0,1)+S(2,2) = 0+0-0+1 = 1$. Let's remember these for future computations: $S(8,6) = S(5,4) = 1$.

 $S(8,5) = S(6,5)+2S(6,4)-S(3,2)+S(5,3) = 0+2S(6,4)-3+7$. $S(6,4) = S(4,4) + 2S(4,3)-S(1,1)+S(3,2) = 0+2 \cdot 2 - 1 + 3 = 6 \implies S(8,5) = 2 \cdot 6 + 4 = 16$. We'll remember $S(6,4) = 6$, $S(8,5) = 16$.

 Lastly $S(8,4) = S(6,4) + 2S(6,3)-S(3,1)+S(5,2) = 6+2 \cdot 16 - 3 + 10 = 45$

 Therefore 
 \[\sum_{i = 0}^{8} S(8,8-i)S(6,i) = S(8,6)S(6,2) + S(8,5)S(6,3) + S(8,4)S(6,4)\]
 \[= 1 \cdot 15 + 16 \cdot 16 + 45 \cdot 6 = 541\]

 
$\textbf{Cases 2 and 3}$

 Rewrite 
 \[\sum_{i = 0}^{8} S(8,8-i) \cdot (S(6, i-1)-S(3, i-3)) + S(6, i-1) \cdot (S(8, 8-i)-S(5, 6-i))\]
 \[= \sum_{i = 0}^{8} 2S(8,8-i)S(6,i-1) - S(8,8-i)S(3,i-3) - S(6,i-1)S(5,6-i)\]
 \[= 2 \sum_{i=2}^{5} S(8,8-i)S(6,i-1) - \sum_{i=3}^{5} S(8,8-i)S(3,i-3) - \sum_{i=2}^{5} S(6,i-1)S(5,6-i)\]

 $\textbf{Case 2 and 3 a.}$

 
 \[\sum_{i=2}^{5} S(8,8-i)S(6,i-1) = S(8,6)S(6,1) + S(8,5)S(6,2) + S(8,4)S(6,3) + S(8,3)S(6,4)\]
 \[= 1 \cdot 6 + 16 \cdot 15 + 45 \cdot 16 + 50 \cdot 6 = 1266\]

 $\textbf{Case 2 and 3 b.}$

 
 \[\sum_{i=3}^{5} S(8,8-i)S(3,i-3) = S(8,5)S(3,0) + S(8,4)S(3,1)+S(8,3)S(3,2)\]
 \[= 16 \cdot 1 + 45 \cdot 3 + 50 \cdot 3 = 301\]

 $\textbf{Case 2 and 3 c.}$

 
 \[\sum_{i=2}^{5} S(6,i-1)S(5,6-i) = S(6,1)S(5,4) + S(6,2)S(5,3) + S(6,3)S(5,2) + S(6,4)S(5,1)\]
 \[= 6 \cdot 1 + 15 \cdot 7 + 16 \cdot 10 + 6 \cdot 5 = 301\]

 Putting it all together, in this case we have

 
 \[\sum_{i = 0}^{8} S(8,8-i)S(6,i-1) = 2 \cdot 1266 - 301 - 301 = 2 \cdot 965 = 1930\]

 $\textbf{Case 4}$

 Finally, 
 \[\sum_{i=0}^{8} S(7,8-i)S(5,i-2) = \sum_{i=3}^{6} S(7,8-i)S(5,i-2)\]
 \[= S(7,5)S(5,1)+S(7,4)S(5,2)+S(7,3)S(5,3)+S(7,2)S(5,4)\]

 Now $S(7,5) = S(5,5)+2S(5,4)-S(2,2)+S(4,3) = 0+2-1+2 = 3$ and $S(7,4) = S(5,4)+2S(5,3)-S(2,1)+S(4,2) = 1+2 \cdot 7 - 2 + 6 = 19$.

 Therefore, 
 \[\sum_{i=0}^{8} S(7,8-i)S(5,i-2) = 3 \cdot 5 + 19 \cdot 10 + 30 \cdot 7 + 21 \cdot 1 = 436\]

 


 
Finally,

 
 \[S(16,8) = \sum_{i = 0}^{8} S(8, 8-i)S(6, i) + S(8,8-i) \cdot (S(6, i-1)-S(3, i-3))\]
 \[+ S(6, i-1) \cdot (S(8, 8-i)-S(5, 6-i)) + S(7,8-i)S(5,i-2)\]

 
 \[= \sum_{i = 2}^{4} S(8,8-i)S(6,i) + (2 \sum_{i=2}^{5} S(8,8-i)S(6,i-1) - \sum_{i=3}^{5} S(8,8-i)S(3,i-3) - \sum_{i=2}^{5} S(6,i-1)S(5,6-i)) + \sum_{i=3}^{6} S(7,8-i)S(5,i-2)\]

 
 \[= 541 + 1930 + 436 = 2907\]

 Thus the answer is 
 \[\boxed{907}\]

 


 





\textbf{Solution 2 and 3 Remark}

There is a clarification for solution 2, but really solution 2 and 3 follow the same two main concepts. Casework and stars and bars. Solution 2 directly solves each case while solution 3 includes each case in its summation.~Bigbrain\_2009

 





\textbf{Solution 2 (Casework)}

Let's split the problem into a few cases:

 Case 1: All $8$ people are sitting isolated (no person sits next to any of them): $^8C_0 \cdot ^9C_1 = 9$

 Case 2: $6$ people are isolated, $2$ people sit next to each other (such that each person sits next to either $0$ or $1$ other person): $^7C_1 \cdot ^9C_2 = 7 \cdot 36 = 252$

 Case 3: $4$ people are isolated, $2$ people sit next to each other and $2$ other people sit next to each other with the $2$ groups of $2$ people not sitting next to each other (so each person still sits next to either $0$ or $1$ other person): $^6C_2 \cdot ^9C_3 = 1260$

 Case 4: $2$ people are isolated, $6$ people are split into $3$ groups of $2$ people, and no $2$ groups sit next to each other: $^5C_3 \cdot ^9C_4 = 10 \cdot 126 = 1260$

 Case 5: $4$ groups of $2$, no groups are sitting next to each other: $^4C_4 \cdot ^9C_5 = 126$

 We have $N = 9 + 252 + 1260 + 1260 + 126 = 2907$

 $2907 \equiv \boxed{907} \pmod{1000}$

 


 \textbf{Clarification for the chooses} For each case there are $8$ chairs taken by the people added with the number of chairs it takes to split the groups of people. For example, with $2$ groups of $2$ people sitting next to each other and $4$ groups of $1$ person, an example will be p p c p p c p c p c p c p. This results in there being $16-8-(2+4-1)=3$ chairs available to distribute (subtracting off the number of people and chairs seperating groups). By stars and bars with the bars being the groups of people, there will be $^9C_{number of remaining chairs}$ ways to distribute the remaining chairs. In this case, $^9C_{3}$. Then, you multiply by the number of ways the groups of two people sitting next to each other and groups of one person can be arranged. This would be the total number of groups choose the number of groups of two people. (To clarify the previous two sentences, you start with $8$, $7$, $6$, $5$, or $4$ groups depending on the case, each group being a single or pair, and you are choosing to make $0$, $1$, $2$, $3$, or $4$ of these groups into pairs instead of singles, respectively.) In this case, $2+4=6$ total groups choose $2$ groups of two people. Then, you multiply these two chooses because both arrangements are independent of each other and sum them over all cases!

 





\textbf{Solution 3 (Stars-and-bars)}

Let $n$ be the number of groups of adjacent people.

 Clearly, there will be $8-n$ groups with $2$ people, and there are $\binom{n}{8-n}$ ways to select these groups.

 Now, we use the \href{https://artofproblemsolving.com/wiki/index.php?title=Ball-and-urn}{stars-and-bars technique} to find the number of ways to distribute the $8$ remaining chairs around the $8$ people. Since the $n$ groups are separated, we must choose $n-1$ chairs to place in between beforehand, leaving $8-(n-1)=9-n$ chairs to distribute among $n+1$ spots. Using stars-and-bars, we find that the number of ways to do this is $\binom{9}{9-n}$.

 Thus, the total number of arrangements is $\binom{n}{8-n} \binom{9}{9-n}$. Summing this from $n=4$ to $n=8$ yields $\sum_{n=4}^{8} \binom{n}{8-n} \binom{9}{9-n} = 2907 \equiv \boxed{907} \pmod{1000}$.

 \href{https://mathproblemsolvingskills.wordpress.com/}{MathProblemSolvingSkills.com}

 


 




	\item (\textit{AIME II 2025}) Let $S$ be the set of vertices of a regular $24$-gon. Find the number of ways to draw $12$ segments of equal lengths so that each vertex in $S$ is an endpoint of exactly one of the $12$ segments.


 
 \textbf{Answer: }113



\textbf{Solution}

The segments we draw must be of equal length, corresponding to a specific step size $k$ (number of steps between vertices).

 For each step size $k$, we need to determine if it is possible to form a perfect matching (non-overlapping segments covering all vertices). The number of such perfect matchings depends on the greatest common divisor (gcd) of $k$ and 24.

 When choosing a step size $k$, the 24-gon is decomposed into $\gcd(k, 24)$ cycles, each of length $\frac{24}{\gcd(k, 24)}$. For a perfect matching to exist, each cycle must be of even length (since each line segment consists of 2 points, the length of the cycle must be divisible by 2).

 For each valid step size ($k$):

 If the cycle length is 2 (diameters), there is exactly 1 way to match the vertices.

 For other even cycle lengths, each cycle contributes a factor of 2 to the number of perfect matchings.

 ($k = 1$): $\gcd(1, 24) = 1$, cycle length 24, 2 matchings.

 ($k = 2$): $\gcd(2, 24) = 2$, cycle length 12, $(2^2 = 4)$ matchings.

 ($k = 3$): $\gcd(3, 24) = 3$, cycle length 8, $(2^3 = 8)$ matchings.

 ($k = 4$): $\gcd(4, 24) = 4$, cycle length 6, $(2^4 = 16)$ matchings.

 ($k = 5$): 

 




	\item (\textit{AIME II 2025}) Let $A_1A_2\dots A_{11}$ be an $11$-sided non-convex simple polygon with the following properties:


 • For every integer $2 \leq i \leq 10$, the area of $A_iA_1A_{i+1}$ is $1$,


 • For every integer $2 \leq i \leq 10$, $\cos(\angle A_iA_1A_{i+1})=\frac{12}{13}$,


 • The perimeter of the $11$-gon $A_1A_2\dots A_{11}$ is equal to $20$.


 Then $A_1A_2+A_1A_{11}$ can be expressed as $\frac{m\sqrt{n}-p}{q}$ where $m,n,p,q$ are positive integers, $n$ is not divisible by any square, and no prime divides all of $m$,$p$, and $q$. Find $m+n+p+q$.


 
 \textbf{Answer: }019



\textbf{Solution 1}

Set $A_1A_2 = x$ and $A_1A_3 = y$. By the first condition, we have $\frac{1}{2}xy\sin\theta = 1$, where $\theta = \angle A_2 A_1 A_3$. Since $\cos\theta = \frac{12}{13}$, we have $\sin\theta = \frac{5}{13}$, so $xy = \frac{26}{5}$. Repeating this process for $\triangle A_i A_1 A_{i+1}$, we get $A_1A_2 = A_1A_4 = \ldots A_1A_{10} = x$ and $A_1A_3 = A_1A_5 = \ldots A_1A_{11} = y$. Since the included angle of these $9$ triangles is $\theta$, the square of the third side is 
 \[x^2 + y^2 - 2xy\cos\theta = x^2 + y^2 - \frac{52}{5}\cdot \frac{12}{13} = x^2 + y^2 - \frac{48}{5} = (x+y)^2 - 20.\]Thus the third side has length $\sqrt{(x+y)^2 - 20}.$ The perimeter is constructed from $9$ of these lengths, plus $A_{11}A_1 + A_1A_2 = x + y$, so $9\sqrt{(x+y)^2 - 20} + x + y = 20$. We seek the value of $x + y,$ so let $x + y = a$ so 
 \begin{align*} 9\sqrt{a^2 - 20} + a &= 20\\ 81(a^2 - 20) &= 400 - 40a + a^2\\ 4a^2 + 2a - 101 &= 0 \\ a &= \frac{-2 \pm \sqrt{1620}}{8} = \frac{-1 \pm \sqrt{405}}{4} = \frac{-1 \pm 9\sqrt{5}}{4}. \end{align*}
Taking the positive solution gives $m + n + p + q = 1 + 9 + 5 + 4 = \boxed{\textbf{(019)}}.$

 -Benedict T (countmath1)

 





\textbf{Solution 2 (Very similar to Solution 1)}

Let $A_1 A_2 = a_1, A_1 A_3 = a_2, ... A_1 A_{10} = a_9, A_1 A_11 a_{10}$. Now the area of each triangle is equal to half the product of two side lengths and the sine of the angle in between. This angle's cosine is given to us to be $\frac{12}{13}$. We recognize this as a $5-12-13$ triangle hence the sine of the angle will be simply $\frac{5}{13}$ (alternatively, one could use the Pythagorean Identity $\sin^2(\angle A_iA_1A_{i+1})+\cos^2(\angle A_iA_1A_{i+1})=1$). Hence each area will be $\frac{1}{2} \cdot a_n a_{n + 1} \cdot \frac{5}{13} = 1 \implies a_n a_{n+ 1} = \frac{26}{5}$. Therefore, we have:

 $a_1 a_2 = \frac{26}{5}$

 $a_2 a_3 = \frac{26}{5}$...

 $a_9 a_{10} = \frac{26}{5}$

 $a_{10} a_1 = \frac{26}{5}$

 Now note that the perimeter is the sum of all the other lengths. For instance, $A_2 A_3$ will be part of this perimeter which happens to be the only unknown side of $\triangle A_1 A_2 A_3$. By Law of Cosines

 $(A_2 A_3)^2 = a_1^{2} + a_2^{2} - 2a_1 a_2 \frac{12}{13}$ after substituting the cosine value. We know $a_1 a_2 = \frac{26}{5}$ so substituting this in we get

 $(A_2 A_3)^2 = a_1^{2} + a_2^{2} - \frac{48}{5} \implies A_2 A_3 = \sqrt{a_1^{2} + a_2^{2} - \frac{48}{5}}$. By symmetry, we conclude

 $\sqrt{a_1^{2} + a_2^{2} - \frac{48}{5}} + ... + \sqrt{a_{10}^{2} + a_1^{2} - \frac{48}{5}} + a_1 + a_2 = 20$ because we need to account for $A_1 A_2 + A_1 A_{11}$ as they are actually known sides that are included in the perimeter.

 Note that from the $a_1 a_2 = \frac{26}{5}$ system of equations, we can see clearly that

 $a_1 = a_3$

 $a_2 = a_4$...

 $a_8 = a_{10}$

 $a_9 = a_1$

 $a_{10} = a_2$

 So we see that $a_{2n + 1} = a_1$ and $a_{2n} = a_2$ so we can substitute this in to get

 $\sqrt{a_1^{2} + a_2^{2} - \frac{48}{5}} + ... + \sqrt{a_1^{2} + a_2^{2} - \frac{48}{5}} + a_1 + a_2 = 20$

 Now the square root terms happen $9$ times as there are $11$ sides and two of them are $a_1, a_2$ given at the end of the $LHS$. So we conclude:

 $9\sqrt{a_1^{2} + a_2^{2} - \frac{48}{5}} + a_1 + a_2 = 20$

 Note that $a_1^{2} + a_2^{2} - \frac{48}{5} = (a_1 + a_2)^2 - 2a_1 a_2 - \frac{48}{5} = (a_1 + a_2)^2 - \frac{52 + 48}{5} = (a_1 + a_2)^2 - 20$. Note that we want to find $A_1 A_2 + A_1 A_{11} = a_1 + a_{10} = a_1 + a_2$. So we let $a_1 + a_2 = x$ to get

 
 \[9\sqrt{x^2 - 20} + x = 20 \implies 81(x^2 - 20) = (20 - x)^2 \implies 81(x^2 - 20) = 400 - 40x + x^2 \implies 81x^2 - 1620 = 400 - 40x + x^2 \implies 80x^2 + 40x - 2020 = 0 \implies 4x^2 + 2x - 101 = 0 \implies x = \frac{-1 \pm 9\sqrt{5}}{4}\]Now take the positive value to get

 $x = \frac{9\sqrt{5} - 1}{4} \implies 9 + 5 + 1 + 4 = \boxed{19}$.

 \textbf{Visualization Tip} The polygon can be imagined sort of like one of those Chinese fans. Congruent triangles branch out from $A_1$, the handle of the fan, forming a series of triangular ridges.

 




	\item (\textit{AIME II 2025}) Let the sequence of rationals $x_1,x_2,\dots$ be defined such that $x_1=\frac{25}{11}$ and

 \[x_{k+1}=\frac{1}{3}\left(x_k+\frac{1}{x_k}-1\right)\] for all $k \geq 1$. Then $x_{2025}$ can be expressed as $\frac{m}{n}$ for relatively prime positive integers $m$ and $n$. Find the remainder when $m+n$ is divided by $1000$.


 
 \textbf{Answer: }248



\textbf{Solution 1 (complete)}

This problem can be split into three parts, listed below:

 \textbf{Part 1: Analyzing Fractions} Let $x_k=\cfrac{a_k}{b_k}$, where $a_k,b_k$ are relatively prime positive integers. First, we analyze the moduli of the problem. Plugging in for $x_2$ yields $x_2=\frac{157}{275}$. Notice that in both $x_1$ and $x_2$, the numerator is equivalent to $1$ and the denominator is equivalent to $2$ modulus $3$. We see that $x_2=\frac{1}{3}\cdot\frac{(a_1-b_1)^2+a_1b_1}{a_1b_1}$. Specifically, we know that
 \[(a_1-b_1)^2+a_1b_1\equiv(1-2)^2+1\cdot2\equiv0\hspace{2mm}(\text{mod}\hspace{1mm}3)\]Then this is always divisible by $3$ for all $x_k$ (it can be shown that for all $x_k$, we have $a_k\equiv1\hspace{2mm}(\text{mod}\hspace{1mm}3)$ and $b_k\equiv2\hspace{2mm}(\text{mod}\hspace{1mm}3)$ by using $\text{mod}\hspace{1mm}9$). Thus, $x_2=\frac{\frac{1}{3}((a_1-b_1)^2+a_1b_1)}{a_1b_1}$, and the numerator and denominator of the right-hand side (RHS) correspond to the numerator and denominator of $x_2$ in simplest form. (To further prove that the top and bottom are relatively prime, consider that $a_k$ and $b_k$ are by definition relatively prime, so $(a_k-b_k)^2$ and $a_kb_k$ share no factors.)

 Notice that the above do not just apply to $x_1$; we did not use any specific properties of $x_1$. Then we may generalize the above, finding that:
 \[a_k=\frac{1}{3}((a_{k-1}-b_{k-1})^2+a_{k-1}b_{k-1})\]
 \[b_k=a_{k-1}b_{k-1}\]Summing the equations yields $a_k+b_k=\frac{1}{3}(a_{k-1}+b_{k-1})^2$ after some manipulation. Let $z_k=a_k+b_k$; then $z_k=\frac{1}{3}z_{k-1}^2$. We are tasked with finding $z_{2025}$.

 \textbf{Part 2: Recursion} We now need an explicit formula for $z_k$. We can first experiment with the recursive formula:

 $z_k=\frac{1}{3}z_{k-1}^2=\frac{1}{3}\left(\frac{1}{3}z_{k-2}^2\right)^2=\frac{1}{3}\left(\frac{1}{3}\left(\frac{1}{3}z_{k-3}^2\right)^2\right)^2$

 Notice that the inner $\frac{1}{3}$ is acted upon by two consecutive powers of $2$. This means that it contributes $\left(\left(\frac{1}{3}\right)^2\right)^2=\left(\frac{1}{3}\right)^4$ to the value of $z_k$. The next innermost $\frac{1}{3}$ is acted upon by one power of $2$, so it contributes $\left(\frac{1}{3}\right)^2$ to the value of $z_k$. Finally, the outermost $\frac{1}{3}$ is acted upon by no powers of $2$, so it contributes $\left(\frac{1}{3}\right)^1$ to the value of $z_k$. The overall power of $\frac{1}{3}$ in $z_k$ in terms of $z_{k-3}$ is then $4+2+1=2^2+2^1+2^0=2^3-1$. Then, the overall power of $\frac{1}{3}$ in $z_k$ in terms of $z_{k-a}$ is $2^a-1$ for positive integers $a$.

 We also see that the $z_{k-3}$ term is acted upon by $3$ powers of $2$, meaning that its power is $2\cdot2\cdot2=2^3$. We can generalize this, so some $z_{k-a}$ term's power is then $2^a$.

 If we combine the above, we obtain the formula $z_k=\left(\frac{1}{3}\right)^{2^a-1}z_{k-a}^{2^a}$. Setting $k=a+1$ results in
 \[z_{a+1}=\left(\frac{1}{3}\right)^{2^a-1}z_1^{2^a}=36^{2^a}\cdot\left(\frac{1}{3}\right)^{2^a-1}\]We can simplify this, noting that $36^{2^a}=12^{2^a}\cdot3^{2^a}$:
 \[z_{a+1}=12^{2^a}\cdot3^{2^a}\cdot\left(\frac{1}{3}\right)^{2^a-1}=3\cdot12^{2^a}\]Finally, decrementing $a+1$ to $a$ gives us our explicit equation:
 \[z_a=3\cdot12^{2^{a-1}}\]

 \textbf{Part 3: Mod Bash} Noting that $z_{2025}=3\cdot12^{2^{2024}}$, we are asked to find its value mod $1000$. We can split mod $1000$ into mod $125$ and mod $8$. We know that $z_{2025}\equiv0\hspace{2mm}(\text{mod}\hspace{1mm}8)$, so we must find its value mod $125$.

 We find that $\phi(125)=100$, so by \href{https://artofproblemsolving.com/wiki/index.php?title=Euler\%27s\_Totient\_Theorem}{Euler's Totient Theorem} we know that $12^{100}\equiv1\hspace{2mm}(\text{mod}\hspace{1mm}125)$. Then, since the power of $12$ is $2^{2024}$, we must find this over modulus $100$.

 Again, we split mod $100$ into mod $4$ and mod $25$. We know that $2^{2024}\equiv0\hspace{2mm}(\text{mod}\hspace{1mm}4)$. Since $\phi(25)=20$, we can apply Euler again, finding that $2^{20}\equiv1\hspace{2mm}(\text{mod}\hspace{1mm}25)$. Then
 \[2^{2024}\equiv(2^{20})^{101}\cdot2^4\equiv2^4\equiv16\hspace{2mm}(\text{mod}\hspace{1mm}25)\]Combining this with the mod $4$ result yields $2^{2024}\equiv16\hspace{2mm}(\text{mod}\hspace{1mm}100)$.

 Going back, $12^{2^{2024}}\equiv12^{16}\hspace{2mm}(\text{mod}\hspace{1mm}125)$. We can then decrement this using a series of simplifications:
 \[12^{16}\equiv144^8\equiv19^8\equiv361^4\equiv(-14)^4\equiv196^2\equiv(-54)^2\equiv2916\equiv41\hspace{2mm}(\text{mod}\hspace{1mm}125)\]Remember that the original value of $z_{2025}$ included a multiplication of $3$; thus,
 \[z_{2025}\equiv41\cdot3\equiv123\hspace{2mm}(\text{mod}\hspace{1mm}125)\]Finally, combining this with the fact that $z_{2025}\equiv0\hspace{2mm}(\text{mod}\hspace{1mm}8)$, we find that the solution to the system of moduli is $z_{2025}\equiv\boxed{248}\hspace{2mm}(\text{mod}\hspace{1mm}1000)$.

 





\textbf{Solution 2 (Faster)}

Note that $x_{k+1} = \frac{1}{3}\left( \frac{(x_k)^{2} - x_k + 1}{x_k} \right)$. An astute reader might recognize the top part as one part of a sum of cubes. I multiplied the entire expression by $x_k + 1$, moved things around a bit, simplified, and was left with the following generalization:$x_{k+1} = \frac{(x_k)^{3} + 1}{3x_k(x_k + 1)}$. Now, we do the following:Set $x_k = \frac{m_k}{n_k}$. Therefore, $x_{k+1} = \frac{m_{k+1}}{n_{k+1}}$. We plug these expressions into the $x_k$ and $x_{k+1}$ and simplify to get: $\frac{m_{k+1}}{n_{k+1}} = \frac{(m_k)^{3} + (n_k)^{3}}{3(m_k)(n_k)(m_k + n_k)}$. Now, as we are looking for the sum of the numerators and denominators of $x_{2025}$, this is great! Now, recall that we want the fraction to be simplest. So we have to cancel out anything we can. Canceling out the factor of $m_k + n_k$ from the numerator and denominator leaves us with $\frac{m_{k+1}}{n_{k+1}} = \frac{(m_k)^{2} - (m_k)(n_k) + (n_k)^{2}}{3(m_k)(n_k)}$. Now, adding the numerator and denominator as well as keeping the extra factor of $3$, we get: 3($m_{k+1} + n_{k+1}) = (m_k)^{2} + 2(m_k)(n_k) + (n_k)^{2}$. Nicely, we get the recursion that $m_{k+1} + n_{k+1} = \frac{(m_k + n_k)^{2}}{3}$. Now, by listing out terms using this recursion and doing mod(1000), we get our answer of $\boxed{248}$.

 
Note: this works, but why are we able to add the numerator and denominator? What if one of the fractions is not fully simplified?

 Answer: Because $m_k$ and $n_k$ are coprime, any prime factor $p \mid m_k$ can not be a factor of $(n_k)^{2}$, any prime factor $q \mid n_k$ can not be a factor of $(m_k)^{2}$. The only possible common factor of the numerator and denominator could be 3. However, a simple induction argument shows $n_k \equiv 2 \mod 3$ and $m_k \equiv 1 \mod 3$ for $k > 1$. So this double recursion $\frac{m_{k+1}}{n_{k+1}} = \frac{\frac{1}{3}\big((m_k)^{2} - (m_k)(n_k) + (n_k)^{2}\big)}{(m_k)(n_k)}$ is indeed the most simplified one.

 Another answer to the same question: Rakko12 has already answered the question with a strong argument. I'm merely just proposing another way to answer it.

 Let $x_k = \frac{m_k}{n_k}$ as a reduced fraction as above. Then, 

 $x_{k + 1} = \frac{m_{k + 1}}{n_{k + 1}}$ and thus

 $\frac{m_{k + 1}}{n_{k + 1}} = \frac{1}{3}(\frac{\frac{m_k^2}{n_k^2} + 1 - \frac{m_k}{n_k}}{\frac{m_k}{n_k}})$ so simplifying yields

 $\frac{m_{k + 1}}{n_{k + 1}} = \frac{1}{3}(\frac{m_k^2 - m_k n_k + n_k^2}{m_k n_k}) = \frac{m_k^2 - m_k n_k + n_k^2}{3m_k n_k}$. Let $m_k = a, n_k = b$. We will find the $\gcd(x^2 - xy + y^2, 3xy)$ when we know that since $x_k$ is a reduced fraction, we have $\gcd(m_k, n_k) = \gcd(x, y) = 1$.

 Note that $\gcd(x^2 - xy + y^2, 3xy) = \gcd(x^2 - xy + y^2 + 3xy, 3xy) = \gcd((x + y)^2, 3xy)$

 Note that $\gcd(x, y) = \gcd(x + y, y) = 1$ and since $x + y, y$ are relatively prime, so are $(x + y)(x + y)$ and $y$ i.e. $\gcd((x + y)^2, y) = 1$. Similarly, from $\gcd(x, y) = \gcd(y, x) = \gcd(y + x, x) = \gcd(x + y, x) = 1 \implies \gcd((x + y)^2, x) = 1$ thus since both of these are equal to $1$, we can safely multiply them to get

 $\gcd((x + y)^2, xy) = 1$. 

 Now if $3 | (x + y)^2 \implies \gcd((x + y)^2, 3xy) = 3$$\gcd((x + y)^2, 3xy) = 1$ otherwise.

 Note that for $3 | (x + y)^2$, it would mean $\frac{m_{k + 1}}{n_{k + 1}} = \frac{m_k^2 - m_k n_k + n_k^2}{3m_k n_k}$ is not fully simplified unless both the numerator and denominator were divided by $3$. Thus, we get

 $\frac{m_{k + 1}}{n_{k + 1}} = \frac{m_k^2 - m_k n_k + n_k^2}{3m_k n_k} = \frac{\frac{m_k^2 - m_k n_k + n_k^2}{3}}{m_k n_k}$ and thus $m_{k + 1} + n_{k + 1} = \frac{m_k^2 - m_k n_k + n_k^2}{3} + m_k n_k = \frac{m_k^2 - m_k n_k + n_k^2 + 3m_k n_k}{3} = \frac{(m_k + n_k)^2}{3}$ but recall that this was only true if $3 | (x + y)^2 = (m_k + n_k)^2$ and thus this would always be an integer. So this proves the claim. A similar argument follows when $3$ does not divide $(x + y)^2$.

 \href{https://mathproblemsolvingskills.wordpress.com/}{MathProblemSolvingSkills.com}

 


 




	\item (\textit{AIME II 2025}) Let ${\triangle ABC}$ be a right triangle with $\angle A = 90^\circ$ and $BC = 38.$ There exist points $K$ and $L$ inside the triangle such
 \[AK = AL = BK = CL = KL = 14.\]The area of the quadrilateral $BKLC$ can be expressed as $n\sqrt3$ for some positive integer $n.$ Find $n.$


 
 \textbf{Answer: }104



\textbf{Solution 1}

From the given condition, we get $\angle{LAK}=60^{\circ}$ and $\triangle{LCA}, \triangle{BAK}$ are isosceles. Denote $\angle{BAK}=\alpha, \angle{CAL}=30^{\circ}-\alpha$. From the isosceles condition, we have $\angle{BKA}=180^{\circ}-2\alpha, \angle{CLA}=120^{\circ}+2\alpha$.

 Since $\angle{CAB}$ is right, then $AB^2+AC^2=BC^2$, we can use the Law of Cosines to express $AC^2$ and $AB^2$, and sum the equations to get $AC^2+AB^2=2\cdot 14^2(2-\cos \angle{BKA}-\angle {CLA})=2\cdot 14^2(2+\cos(2\alpha)+\cos(60^{\circ}-2\alpha))=38^2$.

 This simplifies to $\cos(2\alpha)+\cos(60^{\circ}-2\alpha)=\frac{165}{98}$, and expanding the expression using the angle subtraction formula, we get $\sqrt{3}\sin(2\alpha+60^{\circ})=\frac{165}{98}, \sin(2\alpha+60^{\circ})=\frac{55\sqrt{3}}{98}$.

 Connecting $CK$, we notice that $\angle{CLK}=360^{\circ}-\angle{CLA}-\angle{ALK}=180^{\circ}-2\alpha=\angle{AKB}$, and since $CL=LK=AK=KB$, we have $\triangle{CLK}\cong \triangle{AKB}$ by SAS Congruence. Moreover, since $K$ lies on the perpendicular bisector of $AB$, the distance from $K$ to $AC$ is half of the length of $AB$, which means $[ACK]=\frac{[ABC]}{2}$, so we get $[ACK]=[ACL]+[ALK]+[CLK]=[ACL]+[ALK]+[ABK]=[ABC]-[BKLC]$, which means that $[BKLC]=[AKC]=\frac{[ABC]}{2}$. Thus, we find $[AKC]=[ALK]+\frac{14^2\cdot\sin(\angle{CLA})}{2}+\frac{14^2\cdot\sin(\angle{BKA})}{2}=[ALK]+\frac{14^2}{2}(\sin(60^{\circ}-2\alpha)+\sin 2\alpha)=98(\sin(60^{\circ}+2\alpha))+[ALK]=55\sqrt{3}+\frac{\sqrt{3}}{4}14^2=104\sqrt{3}$. Our answer is $\boxed{104}$.

 





\textbf{Solution 2}


                    \begin{figure}[H]
                        \centering
                        \includegraphics[width=0.5\textwidth]{Solution/AIME_II_2025_14_Solution_Figure_1.png}
                    \end{figure}
                    Let $O$ be the midpoint of $BC$. Take the diagram and rotate it $180^{\circ}$ around $O$ to get the diagram shown. Notice that we have $\angle ABC+\angle ACB=90^{\circ}$. Because $\triangle AKL$ is equilateral, then $\angle KAL=60^{\circ}$, so $\angle BAK+\angle CAL=30^{\circ}$. Because of isosceles triangles $\triangle BAK$ and $\triangle CAL$, we get that $\angle ABK+\angle ACL=30^{\circ}$ too, implying that $\angle KBC+\angle LCB=60^{\circ}$. But by our rotation, we have $\angle LCO=\angle L'BO$, so this implies that $\angle KBL'=60^{\circ}$, or that $\triangle KBL'$ is equilateral (specifically, the presence of a $60^{\circ}$ angle in between two sides of length $14$ allows us to conclude that $\triangle KBL'$ is congruent to an equilateral triangle by SAS). We can similarly derive that $\angle KBO=\angle K'CO$ implies $\angle LCK'=60^{\circ}$ so that $\triangle LK'C$ is also equilateral. At this point, notice that quadrilateral $KL'K'L$ is a rhombus because all its sides have length $14$. The area of our desired region is now $[BKLC]=\frac{1}{2}[BL'K'CLK]$. We can easily find the areas of $\triangle KBL'$ and $\triangle LK'C$ to be $\frac{\sqrt{3}}{4}\cdot 14^2=49\sqrt{3}$ each. Now it remains to find the area of rhombus $KL'K'L$.
                    \begin{figure}[H]
                        \centering
                        \includegraphics[width=0.5\textwidth]{Solution/AIME_II_2025_14_Solution_Figure_2.png}
                    \end{figure}
                    Focus on the quadrilateral $AKOL$. Restate the configuration in another way - we have equilateral triangle $\triangle AKL$ with side length 14, and a point $O$ such that $AO=19$ (the median to the hypotenuse of a right triangle is equal to half the length of the hypotenuse) and $\angle KOL=90^{\circ}$. We are trying to find the area of $\triangle KOL$. Let $M$ be the midpoint of $KL$. We see that $AM=7\sqrt{3}$ ($\overline{AM}$ is a median of equilateral $\triangle AKL$ with side length $14$), and since $M$ is the circumcenter of right $\triangle KOL$, it follows that $MO=7$. Let $\angle KMO=\theta$. From the Law of Cosines in $\triangle AMO$, we can see that 
 \[(7\sqrt{3})^2+7^2-2(7\sqrt{3})(7)\cos (\angle AMO)=361,\] so after simplification we get that $\cos (\theta +90^{\circ})=-\frac{55\sqrt{3}}{98}$. Applying the cosine angle sum identity simplifies this to $\sin \theta =\frac{55\sqrt{3}}{98}$. Then, the Pythagorean Identity $\cos^2\theta +\sin^2\theta =1$ gives us $\cos \theta =\frac{23}{98}$ ($\theta$ is clearly acute, so $\sin \theta$ and $\cos \theta$ are both positive). Applying the Law of Cosines in $\triangle KMO$, we get 
 \[49+49-2\cdot 7\cdot 7\cdot \cos \theta =98-98\cdot \frac{23}{98}=98-23=75=KO^2,\] which tells us $KO=5\sqrt{3}$. The Pythagorean Theorem in $\triangle KOL$ gives that $OL=\sqrt{KL^2-KO^2}={11}$, so the area of $\triangle KOL$ is $\frac{KO\cdot KL}{2}=\frac{55\sqrt{3}}{2}$. The rhombus $KL'K'L$ consists of four of these triangles, so its area is $4\cdot \frac{55\sqrt{3}}{2}=110\sqrt{3}$. 

 Finally, the area of hexagon $BL'K'CLK$ is $49\sqrt{3}+110\sqrt{3}+49\sqrt{3}=208\sqrt{3}$, and since this consists of quadrilaterals $BKLC$ and $CK'L'B$, which must be congruent by rotation, the area of $BKLC$ is $104\sqrt{3}$. Therefore, the answer is $\boxed{104}$.

 





\textbf{Solution 3 (coordinates and bashy algebra)}

By drawing out the triangle, I set $A$ to be $(0, 0)$ in the coordinate plane. I set $C$ to be $(x, 0)$ and B to be $(0, y)$. I set $K$ to be $(a, b)$ and $L$ to be $(c, d)$. Then, since all of these distances are $14$, I used coordinate geometry to set up the following equations:$a^{2} + b^{2} = 196$; $a^{2} + (b - y)^{2} = 196$; $(a - c)^{2} + (b - d)^{2} = 196$; $c^{2} + d^{2} = 196$; $(c - x)^{2} + d^{2} = 196$. Notice by merging the first two equations, the only possible way for it to work is if $b - y = -b$ which means $y = 2b$. Next, since the triangle is right, and we know one leg is $2b$ as $y = 2b$, the other leg, $x$, is $\sqrt{38^{2} - (2b)^{2}}$.We now have:

 $a^{2} + b^{2} = 196  \hspace{1 cm} \textbf{(1)}$

 $(a - c)^{2} + (b - d)^{2} = 196 \hspace{1 cm} \textbf{(2)}$

 $c^{2} + d^{2} = 196  \hspace{1 cm} \textbf{(3)}$

 $(c - \sqrt{38^{2} - (2b)^{2}})^{2} + d^{2} = 196  \hspace{1 cm} \textbf{(4)}$

 Expanding equation (4) and simplifying, we end with $38^{2} - 4b^{2} = 2c \cdot \sqrt{38^{2} - 4b^{2}}$. Next, squaring both sides and canceling terms, we have $4 \cdot 19^{2} - 4b^{2} = 4c^{2}$ which tells us that $c^{2} = 19^{2} - b^{2}$. Now, plugging this value in into equation (3) tells us that $d^{2} = 196 - 19^{2} + b^{2}$. We expand equation (2) to get $a^{2} - 2ac + c^{2} + b^{2} - 2bd + d^{2} = 196$. Using equation (1), we can cancel terms and shift things over to get $196 = 2ac + 2bd$ which means $98 = ac + bd$. From equation (1), we have $a^{2} = 196 - b^{2}$. Now, plugging in all of our variables in terms of $b$ to this new equation, we have $98 = \sqrt{(361 - b^{2})(196 - b^{2})} + b\sqrt{b^{2} - 165}$. We now move things over to get $98 - b\sqrt{b^{2} - 165} = \sqrt{(361 - b^{2})(196 - b^{2})}$. Squaring both sides and canceling, we have $98^{2} - 196b\sqrt{b^{2} - 165} = (361)(196) - 392b^{2}$. We can now divide both sides by $196$ to get $49 - b\sqrt{b^{2} - 165} = 361 - 2b^{2}$. Rearranging and simplifying, we now have $2b^{2} - 312 = b\sqrt{b^{2} - 165}$. Squaring both sides and combining like terms, we have $3b^{4} - 1083b^{2} + 312^{2} = 0$. This part will be a bit of a bash. Quadratic Formula tells us that $b^{2} = \frac{1083 \pm \sqrt{1083^{2} - 12 \cdot 312^{2}}}{6}$. The discriminant nicely simplifies to $\sqrt{4761} = 69$(this will be an extremely long bash but it's worth it). In fact, after computing, we end up with $b^{2} = 192, 169$. This leads us to solutions of $b = 8\sqrt{3}, 13$. If we choose $b = 8\sqrt{3}$, then $c^{2} = 361 - b^{2} = 361 - 192 = 169$ which tells us that $c = 13$.(In fact if you choose $b = 13$, then $c = 8\sqrt{3}$ so it's symmetrical and doesn't matter which one you choose). Next, $a^{2} = 196 - b^{2} = 196 - 192 = 4$ which tells us that $a = 2$. Finally, $d^{2} = 196 - c^{2} = 196 - 169 = 27$ which tells us that $d = 3\sqrt{3}$. Therefore, after all the bashing, our solution quadruple is $a = 2, b = 8\sqrt{3}, c = 13, d = 3\sqrt{3}$. Now plugging in all the points and using the Pythagorean Theorem, we get the coordinates of the quadrilateral which are $(0, 16\sqrt{3}), (2, 8\sqrt{3}), (13, 3\sqrt{3})$, and $(26, 0)$. By Shoelace, our area is $104\sqrt{3}$. Thus, the answer is $\boxed{104}$.

 





\textbf{Solution 4 (Trigonometry)}


                    \begin{figure}[H]
                        \centering
                        \includegraphics[width=0.5\textwidth]{Solution/AIME_II_2025_14_Solution_Figure_3.png}
                    \end{figure}
                    Immediately we should see that $\triangle{AKL}$ is equilateral, so $\angle{KAL}=60$.

 We assume $\angle{LCA}=x$, and it is easily derived that $\angle{KBA}=30-x$. Using trigonometry, we can say that $AC=28\cos{x}$ (imagine extending $\overrightarrow{CL}$ to intersect $\overline{AB}$, forming a right triangle with hypotenuse length $2\cdot14=28$) and $AB=28\cos{(30-x)}$. Pythagoras tells us that $BC^2=AC^2+AB^2$ so now we evaluate as follows:\begin{align*}38^2 &=28^2(\cos^2{x}+\cos^2{(30-x)}) \\(\frac{19}{14})^2 &=\cos^2{x}+(\frac{\sqrt{3}}{2} \cos{x} + \frac{1}{2} \sin{x})^2 \\&=\cos^2{x}+\frac{3}{4} \cos^2{x}+\frac{\sqrt{3}}{2}\sin{x} \cos{x}+\frac{1}{4}\sin^2{x} \\&=\frac{3}{2} \cos^2{x}+\frac{\sqrt{3}}{2}\sin{x} \cos{x}+\frac{1}{4} \\&=\frac{3}{4}(2\cos^2{x}-1)+\frac{\sqrt{3}}{4} (2\sin{x} \cos{x})+1 \\(\frac{33}{14})(\frac{5}{14})&=\frac{\sqrt{3}}{2}(\frac{\sqrt{3}}{2}(\cos{2x})+\frac{1}{2} (\sin{2x})) \\\frac{55\sqrt{3}}{98}&=\cos{(30-2x)} \\\end{align*}


 It is obvious that $\angle{ALC}=180-2x$. We can easily derive $\cos{(150+(30-2x))}$ using angle addition we know, and then the Law of Cosines to find side $AC$.

 \begin{align*}\frac{55\sqrt{3}}{98}=\cos{(30-2x)} \\\sin{(30-2x)}=\sqrt{1-\cos^2{(30-2x)}}=\frac{23}{98} \\\cos{(180-2x)}=(-\frac{\sqrt{3}}{2})(\frac{55\sqrt{3}}{98})-(\frac{1}{2})(\frac{23}{98}) \\\cos{(180-2x)}=-\frac{47}{49} \\AC^2=14^2+14^2-2\cdot 14\cdot 14\cdot (-\frac{47}{49}) \\AC=\sqrt{768}=16\sqrt3 \\\end{align*}


 We easily find $\cos{x}=\frac{4\sqrt{3}}{7}$ and $\sin{x}=\frac{1}{7}$ (draw a perpendicular down from $L$ to $AC$). What we are trying to find is the area of $BKLC$, which can be found by adding the areas of $\triangle{BKL}$ and $\triangle{BLC}$. It is trivial that $\triangle{BKL}$ and $\triangle{ACL}$ are congruent (do some angle chasing in terms of $x$ and use the given side lengths of $14$), so we know that $BL=28\cos{x}$. What we require is

 \begin{align*}\frac{1}{2}(14)(14)(\sin{(180-2x)})+\frac{1}{2}(14)(28\cos{x})(\sin{(120+x)}) \\\end{align*}


 We do similar calculations to obtain that $\sin{(120+x)}=\frac{11}{14}$ and $\cos{(180-2x)}=-\frac{47}{49}$ implies $\sin{(180-2x)}=\frac{8\sqrt{3}}{49}$, so now we plug in everything we know to calculate the area of the quadrilateral:

 \begin{align*}&  \frac{1}{2}(14)(14)(\sin{(180-2x)})+\frac{1}{2}(14)(28\cos{x})(\sin{(120+x)}) \\&=\frac{1}{2}(14)(14)(\frac{8\sqrt{3}}{49})+\frac{1}{2}(14)(16\sqrt{3})(\frac{11}{14}) \\&=16\sqrt{3}+88\sqrt{3} \\&=104\sqrt{3} \\\end{align*}


 We see that $n=\boxed{104}$.

 





\textbf{Solution 5 (Circles and Trigonometry)}


                    \begin{figure}[H]
                        \centering
                        \includegraphics[width=0.5\textwidth]{Solution/AIME_II_2025_14_Solution_Figure_4.PNG}
                    \end{figure}
                    

 Since $KB=KL=KA=14$ and $LK=LA=LC=14$, we can construct 2 circles of radus 14 with $K$ and $L$ as the center of the two circles. Let the intersection of the 2 circles other than $A$ be point $M$. Connect $BM$, $CM$, $KM$, and $LM$. Connect $AM$, which is the radical axis of the 2 circles.

 From the figure, we know that
 \[[KLCB] = [KLCMB] - [BMC]\]
 \[[KLCB] = [BKM] + [CLM] + [KLM] - [BMC]\]

 Let $\angle{BAM} = \theta$, which means that $\angle{CAM} = \frac{\pi}{2} - \theta$. For easier calculation, we temporarily define the radius of the 2 circles (which is 14) to be $R$. $\angle{BAM}$ is an inscribed angle and $\angle{BKM}$ is a central angle, so $\angle{BKM} = 2\angle{BAM} = 2\theta$. Similar with the other side, $\angle{CLM} = \pi-2\theta$. $KM = KL = LM = R$, so $\triangle{KLM}$ is an equilateral triangle.

 Using the Law of Cosines, we get the area of each little triangle, and applying the Law of Sines gives $BM$ and $MC$.
 \[[BKM] = \frac{1}{2}\cdot R^2\cdot\sin(2\theta)\]
 \[[CLM] = \frac{1}{2}\cdot R^2\cdot\sin(\pi-2\theta) = \frac{1}{2}\cdot R^2\cdot\sin(2\theta)\]
 \[[KLM] = \frac{1}{2}\cdot R^{2} \cdot\sin({\frac{\pi}{3}})=\frac{\sqrt3}{4}R^2\]\begin{align*}[BMC] & = \frac{1}{2}\cdot|BM|\cdot|MC|\cdot\sin({\frac{5\pi}{6}})\\&= \frac{1}{2}\cdot\frac{1}{2}\cdot2R\sin(\theta)\cdot2R\sin(\frac{\pi}{2}-\theta)\\&= R^2\cdot\sin(\theta)\cos(\theta)\\&= \frac{1}{2}\cdot R^2\sin(2\theta)\\\end{align*}


 We can conclude that
 \[[KLCB] = \frac{1}{2}\cdot R^2\cdot\sin(2\theta)+\frac{1}{2}\cdot R^2\cdot\sin(2\theta)+\frac{\sqrt3}{4}R^2-\frac{1}{2}\cdot R^2\sin(2\theta)\]
 \[[KLCB] = {14}^2\cdot(\frac{\sin(2\theta)}{2}+\frac{\sqrt3}{4})\]

 Now, we just needed to find the value of $\sin(2\theta)$. We analyze the $\triangle{BMC}$. We already know that $\angle{BMC} = {150}^{\circ}$, $BM = 2R\sin(\theta)$, and $CM = 2R\cos(\theta)$. Using the Laws of Cosines (again!) and the given condition of $BC = 38$, we can create a formula on $\theta$.

 
 \[{BC}^2 = {BM}^2+{CM}^2-2\cdot BM\cdot MC\cdot\cos(\angle{BMC})\]
 \[{BC}^2 = (2R\sin(\theta))^2+(2R\cos(\theta))^2-2\cdot\cos({150}^{\circ})\cdot(2R\cos(\theta))\cdot(2R\cos(\theta)) = {38}^2\]
 \[4R^2(\sin^2(\theta)+\cos^2(\theta)+\sqrt3\sin(\theta)\cos(\theta)) = {38}^2\]
 \[4R^2(1+\frac{\sqrt3}{2}\cdot\sin(2\theta)) = {38}^2\]
 \[4R^2+\frac{4R^2\sqrt3}{2}\cdot\sin(2\theta) = {38}^2\]
 \[\sin(2\theta) = \frac{2}{\sqrt3}\cdot(\frac{{38}^2}{4\cdot{R}^2}-1)\]
 \[\sin(2\theta) = \frac{2\cdot165}{\sqrt3\cdot{14}^2} = \frac{165}{98\sqrt3}\]

 We put the calculated value of $\sin(2\theta)$ back into $[KLCB]$:
 \[[KLCB] = {14}^2\cdot(\frac{165}{2\cdot98\sqrt3}+\frac{\sqrt3}{4})\]
 \[[KLCB] = 55\sqrt3+49\sqrt3 = 104\sqrt3\]

 Therefore, $n=\boxed{104}$.

 





\textbf{Solution 6 (Trig Identities; warning: bashy)}

Consider a diagram to the original problem (credit to solution 4):

 
                    \begin{figure}[H]
                        \centering
                        \includegraphics[width=0.5\textwidth]{Solution/AIME_II_2025_14_Solution_Figure_5.png}
                    \end{figure}
                    

 Now, let us simplify the problem further. We know that $K$ and $L$ must lie on the perpendicular bisectors of $AB$ and $AC$, respectively. The real problem here is the equilateral triangle in the middle, inscribed in a rectangle with diagonal length 19.

 We create a further simplified problem: given that the inscribed equilateral triangle of a certain rectangle with diagonal length $19$ has side length $14$, find the sides and intersection points on this rectangle. For reference, here is a diagram:

 
                    \begin{figure}[H]
                        \centering
                        \includegraphics[width=0.5\textwidth]{Solution/AIME_II_2025_14_Solution_Figure_6.png}
                    \end{figure}
                    

 


 Note the angles $\angle{LAD}$ and $\angle{BAK}$. Since $\angle{LAD} + \angle{BAK} + 60^{\circ} = 90^{\circ}$, $\angle{LAD} + \angle{BAK} = 30^{\circ}$, and $\angle{BAK} = 30^{\circ} - \angle{LAD}$. Thus, let $\angle{LAD} = \alpha$ and $\angle{BAK} = 30 - \alpha$. 

 Now, we know that $AB^2 + AD^2 = 19^2$, as the hypotenuse of the larger right triangle is $38$. However, we can also express $AB$ and $AD$ in terms of $\alpha$: $AB = 14(\cos(30^{\circ}-\alpha))$ and $AD = 14(\cos(\alpha))$. Thus, $\cos^2(\alpha) + \cos^2(30^{\circ}-\alpha) = 361/196$. We expand this using the cosine difference identity:

 $\cos^2(\alpha) + (\cos(30^{\circ})\cos(\alpha) + \sin(30^{\circ})\sin(\alpha))^2 = \frac{361}{196}$

 $\frac{7}{4}\cos^2(\alpha) + \frac{1}{4}\sin^2(\alpha) + \frac{\sqrt3}{2}\sin(\alpha)\cos(\alpha) = \frac{361}{196}$

 Using the fact that $\sin^2(\alpha) + \cos^2(\alpha) = 1$, then multiplying the entire equation by $2$,

 $3\cos^2(\alpha) + \sqrt3\sin(\alpha)\cos(\alpha) = \frac{156}{49}$

 Now, to save some writing, let us denote $\sin(\alpha)$ with $x$, and $\cos(\alpha)$ with $y$.

 We have the following equations:

 $x^2 + y^2 = 1$

 $3y^2 + \sqrt3xy = \frac{156}{49}$

 Substituting $x$ for $y$, moving $3y^2$ to the left side, squaring, and dividing by 9, we end up with the quartic:

 $\frac{4}{3}y^4 - \frac{361}{147}y^2 + \frac{52^2}{49^2} = 0$

 Using the quadratic formula, we end up with this:

 $y^2 = \frac{\frac{361}{49} \pm \frac{1}{49}\cdot\sqrt{361^2 - 208^2\cdot3}}{8}$

 Now, we could just compute $361^2 - 208^2\cdot3$, but instead, we can do this:

 $361^2 - 208^2\cdot3 = (129600 + 720 + 1) - (40000 + 3200 + 64)\cdot3$

 $(129600 + 721) - (43200 + 64)\cdot3$

 $(129600 + 721) - (129600 + 192) = 529 = 23^2$

 Thus, we have two cases:

 $1. \cos(\alpha) = \frac{13}{14}$

 $2. \cos(\alpha) = \frac{4\sqrt3}{7}$

 Both lead to the same side lengths of the rectangle: $8\sqrt3$, and $13$. Referring back to our original rectangle diagram and plugging in our trigonometric values, we get that $CK = 13 - 2 = 11$, and $CL = 8\sqrt3 - 3\sqrt3 = 5\sqrt3$. Thus, the area of the original quadrilateral is $\frac{88\sqrt3 + 55\sqrt3 + 65\sqrt3}{2}$, or $\boxed{104}\sqrt3$.

 





\textbf{Solution 7 (analytic geometry with roots of unity)}


                    \begin{figure}[H]
                        \centering
                        \includegraphics[width=0.5\textwidth]{Solution/AIME_II_2025_14_Solution_Figure_7.png}
                    \end{figure}
                    This diagram is modified from the solution 4 diagram. Let $M$ be the midpoint of $AC$, and let $N$ be the midpoint of $AB$.

 We place the diagram onto the Cartesian coordinate grid. Let $A = (0, 0)$, $M = (a, 0)$, $C = (2a, 0)$, $N = (0, b)$, and $B = (0, 2b)$. We are given $AL = CL$, so $\triangle ACL$ is isosceles. Therefore, $LM$ is the perpendicular bisector of $AC$, so we can let $L = (a, y)$. Similarly, we’re given $AK = BK$, so $\triangle ABK$ is also isosceles, and $NK$ is the perpendicular bisector of $AB$. Therefore, we can let $K = (x, b)$.

 We have $AB = 2b$ and $AC = 2a$. We’re given that $\angle BAC = 90^\circ$ and $BC = 38$, so by the Pythagorean theorem, 
 \[(2a)^2 + (2b)^2 = 38^2 \implies 4a^2 + 4b^2 = 1444 \implies a^2 + b^2 = 361.\]

 We now place the diagram onto the complex plane. We use the x-axis of the coordinate plane as the complex plane’s real axis, and we use the y-axis of the coordinate plane as the complex plane’s imaginary axis. So, on the complex plane, $A = 0$, $L = a + yi$, and $K = x + bi$. Also, since we are given $AK = KL = AL$, $\triangle AKL$ is equilateral. In addition, since $AL = AK$, $\angle KAL = 60^\circ$, and because we constructed our diagram with $K$ counterclockwise of $L$ (if it were the other way around, we could go through the same steps as this solution, but with variables switched around), $K$ is a $60^\circ$ counterclockwise rotation of $L$ about $A$, and $L$ is a $60^\circ$ clockwise or $300^\circ$ counterclockwise rotation of $K$ about $A$.

 Rotation on the complex plane is equivalent to multiplying by a root of unity. Here, $K$ and $L$ are rotated a multiple of $60^\circ$ to each other about $A$. $60^\circ$ is one-sixth of a full circle, so to go from $L$ to $K$ or $K$ to $L$, we multiply by a 6th root of unity. Specifically, to go from $L$ to $K$, we multiply by $\dfrac{1}{2} + \dfrac{\sqrt{3}}{2}i$, and to go from $K$ to $L$, we multiply by $\dfrac{1}{2} - \dfrac{\sqrt{3}}{2}i$.

 We multiply the coordinate of $L$ by $\dfrac{1}{2} + \dfrac{\sqrt{3}}{2}i$ on the complex plane to obtain equations for the coordinates of $K$: 
 \[(a + yi)\left(\dfrac{1}{2} + \dfrac{\sqrt{3}}{2}i\right) = \left(\dfrac{a}{2} - \dfrac{\sqrt{3}}{2}y\right) + \left(\dfrac{y}{2} + \dfrac{\sqrt{3}}{2}a\right)i = x + bi.\] Equating real and imaginary parts, we obtain 
 \[x = \dfrac{a}{2} - \dfrac{\sqrt{3}}{2}y \text{ and } b = \dfrac{y}{2} + \dfrac{\sqrt{3}}{2}a.\]

 Similarly, we multiply the coordinate of $K$ by $\dfrac{1}{2} - \dfrac{\sqrt{3}}{2}i$ to obtain equations for the coordinates of $L$: 
 \[(x + bi)\left(\dfrac{1}{2} - \dfrac{\sqrt{3}}{2}i\right) = \left(\dfrac{x}{2} + \dfrac{\sqrt{3}}{2}b\right) + \left(\dfrac{b}{2} - \dfrac{\sqrt{3}}{2}x\right) = a + yi.\] Equating real and imaginary parts, we obtain 
 \[a = \dfrac{x}{2} + \dfrac{\sqrt{3}}{2}b \text{ and } y = \dfrac{b}{2} - \dfrac{\sqrt{3}}{2}x.\]

 We now look back at the problem to see what it asks for: $[BKLC]$. Looking at the diagram, we see we can express the area of the quadrilateral as the area of the big right triangle $ABC$ minus the two isosceles triangles $ABK$ and $ALC$ minus the equilateral triangle $AKL$: 
 \[[BKLC] = [ABC] - [ABK] - [ACL] - [AKL].\]

 We are given that $AK = KL = AL = 14$, so the area of equilateral triangle $AKL$ is $\dfrac{\sqrt{3}}{4} \cdot 14^2 = 49\sqrt{3}$. Also, we can use $AC = 2a$ as the base of $\triangle ABC$ and $AB = 2b$ as the height, so $[ABC] = \dfrac{(2a)(2b)}{2} = 2ab$. Similarly, we use $AC = 2a$ as the base of $\triangle ACL$ and $ML = y$ ($M = (a, 0)$ and $L = (a, y)$, so the distance between the two is equal to $y$) as the height, so $[ACL] = \dfrac{(2a)(y)}{2} = ay$. Finally, we use $AB = 2b$ and $KN = x$ ($N = (0, b)$ and $K = (x, b)$, so the distance between the two is equal to $x$) as the base and height of $\triangle ABK$ respectively, so $[ABK] = \dfrac{(2b)(x)}{2} = bx$. Therefore, 
 \[[BKLC] = [ABC] - [ABK] - [ACL] - [AKL] = 2ab - bx - ay - 49\sqrt{3}.\]

 We have already shown that $b = \dfrac{y}{2} + \dfrac{\sqrt{3}}{2}a$. Substituting this into $a^2 + b^2 = 361$, we have 
 \[a^2 + \left(\dfrac{y}{2} + \dfrac{\sqrt{3}}{2}a\right)^2 = 361.\] Expanding this out, we have 
 \[a^2 + \dfrac{y^2}{4} + \dfrac{\sqrt{3}}{2}ay + \dfrac{3}{4}a^2 = 361.\] Multiplying both sides by $4$ and rearranging the left side, we have 
 \[7a^2 + y^2 + 2ay\sqrt{3} = 1444.\] We previously showed that $AC \perp ML$, so $AM \perp ML$ (since $M$ is on $AC$). Therefore, $\triangle AML$ has a right angle at $M$. By the Pythagorean Theorem, $AM^2 + ML^2 = a^2 + y^2 = AL^2 = 196$. Subtracting $a^2 + y^2$ from the left side and $196$ from the right side, we obtain 
 \[6a^2 + 2ay\sqrt{3} = 1248.\] Dividing both sides of the equation by $2\sqrt{3}$ and factoring $a$ out of the left side, we have 
 \[a(a\sqrt{3} + y) = 208\sqrt{3}.\] However, we have $b = \dfrac{y}{2} + \dfrac{\sqrt{3}}{2}a$, so the expression inside the parentheses is simply $2b$! Therefore, 
 \[2ab = 208\sqrt{3}.\]

 The algebra’s not over yet. We also showed that $a = \dfrac{x}{2} + \dfrac{\sqrt{3}}{2}b$, so substituting that into $a^2 + b^2 = 361$, we obtain 
 \[b^2 + \left(\dfrac{x}{2} + \dfrac{\sqrt{3}}{2}b\right)^2 = 361.\] Expanding this out, we have 
 \[b^2 + \dfrac{x^2}{4} + \dfrac{\sqrt{3}}{2}bx + \dfrac{3}{4}b^2 = 361.\] Multiplying both sides by $4$ and rearranging the left side, we now have 
 \[7b^2 + x^2 + 2xb\sqrt{3} = 1444.\] Does this equation look familiar? We previously showed that $NK \perp AB$. Therefore, $NK \perp AN$ (since $N$ is on $AB$). So, $\triangle ANK$ has a right angle at $N$. By the Pythagorean Theorem, $AN^2 + KN^2 = b^2 + x^2 = AK^2 = 196$. Subtracting $b^2 + x^2$ from the left side and $196$ from the right side, we have 
 \[6b^2 + 2xb\sqrt{3} = 1248.\] We previously also had the equation $6a^2 + 2ay\sqrt{3} = 1248$, and adding this equation to the above equation and factoring out $2\sqrt{3}$, we have 
 \[6a^2 + 6b^2 + 2\sqrt{3}(bx + ay) = 2496.\] We previously showed $a^2 + b^2 = 361$, so $6a^2 + 6b^2 = 6 \cdot 361 = 2166$. Subtracting $6a^2 + 6b^2$ from the left side and $2166$ from the right side, we obtain $2\sqrt{3}(bx + ay) = 330$. Finally, dividing both sides by $2\sqrt{3}$, we have 
 \[bx + ay = 55\sqrt{3}.\]

 We previously arrived at this expression for $[BKLC]$: 
 \[[BKLC] = 2ab - bx - ay - 49\sqrt{3}.\] We now know $2ab = 208\sqrt{3}$ and $bx + ay = 55\sqrt{3}$, so we can simply substitute them in. Therefore, 
 \[[BKLC] = 2ab - bx - ay - 49\sqrt{3} = 208\sqrt{3} - 55\sqrt{3} - 49\sqrt{3} = 104\sqrt{3}.\] Finally, we are given $[BKLC] = n\sqrt{3}$ for some integer $n$. We know $[BKLC] = 104\sqrt{3}$, so $n = \boxed{104}$.

 Notice that $[BKLC] = 104\sqrt{3} = \dfrac{208\sqrt{3}}{2} = \dfrac{[ABC]}{2}$. Is this a coincidence?

 





\textbf{Solution 8}


                    \begin{figure}[H]
                        \centering
                        \includegraphics[width=0.5\textwidth]{Solution/AIME_II_2025_14_Solution_Figure_8.png}
                    \end{figure}
                    $AK = BK = KL, \angle AKL = 60 ^\circ \implies \angle ABL = 30^\circ$ because an inscribed angle has half the measure of its corresponding central angle.

 Similarly, $\angle ACK = 30^\circ.$
 \[\angle BAK = \angle ABK, \angle LBK + \angle ABK = 30^\circ.\]
 \[\angle CAL + \angle BAK = 90^\circ - \angle LAK = 30^\circ \implies\]
 \[\angle CAL = \angle LBK \implies \triangle KLB = \triangle LAC \implies\]$AC = BL.$ Similarly, $AB = CK.$

 Denote the area of triangle $X$ as $[X]$ ($\leftarrow$ Wait, what's this? ~ Christian).
 \[2[ABL] = AB \cdot BL \cdot \sin 30^\circ = \frac{AB \cdot AC}{2} = [ABC].\]Similarly, $[ABL] = [ACK] = \frac {[ABC]}{2}.$
 \[[BKLC] = [ABC] - [ABL] - [ALC]+ [BLK] = [ABC] - [ABL] = [ABL].\] By applying the Law of Cosines on $\triangle ABL,$  we get$AB^2 + BL^2 - 2 AB \cdot BL \cos 30^\circ = AL^2 \implies  AB^2 + AC^2 - AB \cdot AC \sqrt{3} = AL^2 \implies$$BC^2-2[ABC]\sqrt{3}=AL^2 \implies$
 \[BC^2 - AL^2 = 4 [BKLC] \sqrt{3} \implies\]
 \[[BLKC] = \frac {38^2-14^2}{4 \sqrt{3}} = \frac{19^2-7^2}{3} \sqrt{3} = 4 \cdot 26 \sqrt{3} = 104 \sqrt{3}.\]\textbf{vladimir.shelomovskii@gmail.com, vvsss}

 





\textbf{Solution 9}

$\triangle KAL$ is clearly equilateral, and $\triangle AKB$ and $\triangle ALC$ are clearly isosceles. Now, we can do a bit of angle chasing. Say $\angle AKB = 2\theta$. Then, since $\triangle AKB$ is isosceles, $\angle ABK = 90^\circ - \theta$. $\angle KAL = 60^\circ$, so $\angle LAC = \theta - 60^\circ$, so $\angle ALC = 300^\circ - 2\theta$, so $\angle KLC = 2\theta$. Since the angles are congruent, $\triangle AKB$ is congruent to $\triangle KLC$ by SAS. As such, we can rotate it around the center of $\triangle KAL$ onto side $\overline{AL}$ to produce $\triangle LAD$. Because all the triangles are congruent, connecting the congruent corners will all produce the same length. From $BC$, we see that this same length is $38$, so $\triangle BCD$ is equilateral with side length $38$. This triangle is comprised of 3 quadrilaterals and $\triangle KAL$ in the center, and it's easy to see that each quadrilateral is congruent to $BKLC$. As such, $[BKLC] = \frac{[BCD]-[KAL]}{3} = \frac{361 \sqrt{3} - 49 \sqrt{3}}{3} = 104 \sqrt{3}$, giving an answer of $\boxed{104}$. 

 





\textbf{Solution 10 (Algebraic bashing with some trig)}


                    \begin{figure}[H]
                        \centering
                        \includegraphics[width=0.5\textwidth]{Solution/AIME_II_2025_14_Solution_Figure_9.png}
                    \end{figure}
                    We observe that the desired area can be found by subtracting the areas of the three smaller triangles $\triangle BAK$, $\triangle KAL$, and $\triangle LAC$ from the area of the large triangle $\triangle ABC$. This can be expressed as \([ABC]-[BAK]-[KAL]-[LAC]\). Thus, our goal is to compute these four areas.

 Since $\triangle KAL$ is equilateral with side length $14$, its area is
 \[[KAL]=\frac{14^2\sqrt{3}}{4}=49\sqrt{3}.\]

 Let $AC=a$ and $AB=b$. Because $\angle CAB$ is a right angle, by the Pythagorean Theorem we have
 \[a^2+b^2=1444.\]

 Let $\angle CAL=\alpha$ and $\angle KAB=\beta$. Then
 \[\cos\alpha=\frac{a}{28}\] (imagine extending $\overrightarrow{CL}$ to intersect $\overline{AB}$, forming a right triangle with hypotenuse length $2\cdot14=28$), 
 \[\quad \cos\beta=\frac{b}{28}\] similarly.Using Pythagorean identities,
 \[\sin\alpha=\frac{\sqrt{784-a^2}}{28}, \quad \sin\beta=\frac{\sqrt{784-b^2}}{28}.\]

 Since $\angle CAB=90^\circ$ and $\triangle KAL$ is equilateral, we have
 \[\alpha+\beta=30^\circ.\]Thus,
 \[\cos(\alpha+\beta)=\cos\alpha\cos\beta-\sin\alpha\sin\beta=\frac{\sqrt{3}}{2}.\]

 Substituting,
 \[\frac{\sqrt{3}}{2}=\frac{ab}{784}-\frac{\sqrt{(784-a^2)(784-b^2)}}{784}.\]Rearranging,
 \[\frac{ab}{784}-\frac{\sqrt{3}}{2}=\frac{\sqrt{(784-a^2)(784-b^2)}}{784}.\]Clearing denominators,
 \[ab-392\sqrt{3}=\sqrt{(784-a^2)(784-b^2)}\]Squaring both sides gives
 \[a^2b^2-784ab\sqrt{3}+392^2\cdot3=(784-a^2)(784-b^2).\]Expanding the RHS,
 \[a^2b^2-784ab\sqrt{3}+392^2\cdot3=784^2-784a^2-784b^2+a^2b^2.\]Canceling $a^2b^2$ and using $a^2+b^2=1444$,
 \[-784ab\sqrt{3}+392^2\cdot3=784^2-784(1444).\]Simplifying,
 \[ab\sqrt{3}=1248 \quad \Rightarrow \quad ab=416\sqrt{3}.\]

 The area of the large triangle is
 \[[ABC]=\frac{ab}{2}=208\sqrt{3}.\]

 Now substitute $b=\frac{416\sqrt{3}}{a}$ into $a^2+b^2=1444$:
 \[a^2+\left(\frac{416\sqrt{3}}{a}\right)^2=1444,\]
 \[a^4-1444a^2+416^2\cdot3=0.\]Solving for $a^2$,
 \[a^2=\frac{1444\pm\sqrt{1444^2-12\cdot416^2}}{2}=\frac{1444\pm184}{2},\]so
 \[a^2=768 \text{ or } 676.\]

 From the diagram, $a$ is the longer side, so
 \[a=16\sqrt{3}, \quad b=26.\]

 Points $L$ and $K$ lie on the perpendicular bisectors of $AC$ and $AB$, respectively. Let the heights from $L$ to $AC$ and from $K$ to $AB$ be $h_L$ and $h_K$. Then,
 \[(8\sqrt{3})^2+h_L^2=14^2, \quad 13^2+h_K^2=14^2,\]which gives
 \[h_L=2, \quad h_K=3\sqrt{3}.\]

 Thus,
 \[[LAC]=\frac{16\sqrt{3}\cdot2}{2}=16\sqrt{3},\]
 \[[BAK]=\frac{26\cdot3\sqrt{3}}{2}=39\sqrt{3}.\]

 Finally, subtracting areas,
 \[[ABC]-[BAK]-[KAL]-[LAC]=208\sqrt{3} - 39\sqrt{3} - 49\sqrt{3} - 16\sqrt{3} = 104\sqrt{3}.\]

 Therefore, the answer is $\boxed{104}$.

 





\textbf{Solution 12}

Suppose the coordinates of $A(0, 0)$, $B(0, 2b)$, $C(2a, 0)$, $K(d, b)$, $L(a, c)$ because $K$ and $L$ lie on the perpendicular bisector of $AB$ and $AC$ respectively. According to the lengths given, we have

 \begin{align}(2a)^2 + (2b)^2 &= 38^2 \implies a^2 + b^2 = 19^2 \tag{1}\\a^2 + c^2 &= 14^2 \tag{2}\\b^2 + d^2 &= 14^2 \tag{3}\\(a - d)^2 + (b - c)^2 &= 14^2 \tag{4}\end{align}

 Use shoelace formula, \begin{align}S_{AKL} &= \dfrac{1}{2}\begin{vmatrix}0 & a & d\\0 & c & b\end{vmatrix} = \dfrac{1}{2}(ab - cd) = \dfrac{\sqrt{3}}{4} \times 14^2 = 49\sqrt{3}\\S_{CBKL} &= \dfrac{1}{2}\begin{vmatrix}2a & 0 & d & a\\0 & 2b & b & c\end{vmatrix} = \left|\dfrac{3}{2}ab - bd - ac + \dfrac{1}{2}cd\right|\\&= \left|\dfrac{1}{2}(ab - cd) + (a - d)(b - c)\right|\\&= |49\sqrt{3} + (a - d)(b - c)|\\&= |49\sqrt{3} + xy|\end{align}

 Here, let $x = a - d$, $y = b - c$, $x^2 + y^2 = 14^2$, the goal is to find $xy$. While using elimination with $(1) - (3)$, we get $a^2 - d^2 = 19^2 - 14^2 \implies (a - d)(a + d) = 165 \implies a + d = \dfrac{165}{x} \implies a = \dfrac{x + \dfrac{165}{x}}{2}$. Similarly, $b = \dfrac{y + \dfrac{165}{y}}{2}$.

 We substitute them in equation $(1)$,\begin{align}\left(x + \dfrac{165}{x}\right)^2 + \left(y + \dfrac{165}{y}\right)^2 &= 38^2\\x^2 + y^2 + 660 + \dfrac{165^2(x^2 + y^2)}{x^2y^2} &= 38^2\\14^2 + 660 + \dfrac{165^2 \cdot 14^2}{x^2y^2} &= 38^2\\\dfrac{165^2 \cdot 14^2}{x^2y^2} &= 588 = 3 \cdot 14^2\\xy &= 55\sqrt{3}\end{align}

 Then, 
 \[S_{CBKL} = |49\sqrt{3} + 55\sqrt{3}| = 104\sqrt{3}\]

 The answer is $\boxed{104}$.

 \textbf{Remarks} This problem can be approached either by analytic geometry or by trigonometric manipulation. The characteristics of this problem make it highly similar to \href{https://artofproblemsolving.com/wiki/index.php?title=2017\_AIME\_I\_Problems/Problem\_15}{2017 AIME I Problem 15}.

 




	\item (\textit{AIME II 2025}) There are exactly three positive real numbers $k$ such that the function

 \[f(x)=\frac{(x-18)(x-72)(x-98)(x-k)}{x}\] defined over the positive real numbers achieves its minimum value at exactly two positive real numbers $x$. Find the sum of these three values of $k$.


 
 \textbf{Answer: }240



\textbf{Solution 1 ('clunky', trial and error)}

Let $n$ be the minimum value of the expression (changes based on the value of $k$, however is a constant). Therefore, we can say that\begin{align*}f(x)-n=\frac{(x-\alpha)^2(x-\beta)^2}{x}\end{align*}
This can be done because $n$ is a constant, and for the equation to be true for all $x$, the right side is also a quartic. The roots must also both be double, or else there is an even smaller value, a contradiction.

 We expand as follows, comparing coefficients:

 \begin{align*}(x-18)(x-72)(x-98)(x-k)-nx=(x-\alpha)^2(x-\beta)^2 \\-2\alpha-2\beta=-18-72-98-k \implies \alpha+\beta=94+\frac{k}{2} \\\alpha^2+4\alpha \beta +\beta^2=(-18\cdot -72)+(-18\cdot-98)+(-18\cdot-k)+(-72\cdot-98)+(-72\cdot-k)+(-98\cdot-k)=10116+188k \\(\alpha^2)(\beta^2)=(-18)(-72)(-98)(-k) \implies \alpha \beta=252\sqrt{2k} \\\end{align*}


 Recall $(\alpha+\beta)^2+2\alpha \beta=\alpha^2+4\alpha \beta +\beta^2$, so we can equate and evaluate as follows:

 \begin{align}(94+\frac{k}{2})^2+504\sqrt{2k}=10116+188k \tag{1}\\\end{align}\begin{align*}(47-\frac{k}{4})^2+126\sqrt{2k}=2529 \\\frac{k^2}{16}-\frac{47}{2}k+126\sqrt{2k}-320=0 \\\end{align*}


 We now have a quartic with respect to $\sqrt{k}$. Keeping in mind it is much easier to guess the roots of a polynomial with integer coefficients, we set $a=\frac{k}{8}$. Now our equation becomes

 \begin{align*}4a^2-188a+504\sqrt{a}-320=0 \\a^2-47a+126\sqrt{a}-80=0 \\\end{align*}


 We can test for easy roots and find $\sqrt{a}=1$ and $2$. After this, solving the resulting quadratic gets you the remaining roots as $5$ and $-8$. Discard $-8$. Working back through our substitution for $a$, we have generated values of $k$ as $(8, 32, 200)$.

 The sum of all $k$ then must be $8+32+200=\boxed{240}$.

 Note:I'm not sure what the first author meant by, "or else there is an even more 'minimum' value." The most noticeable reason there are two double roots is that there are two distinct positive solutions per the conditions of the problem

 Further notes (from the author):To clarify why the equation takes the form above, if you do not regard the context of the problem, you could say that $\frac{(x-\alpha)^3(x-\beta)}{x}$ is another form, also having two roots. However when we consider this form, (WLOG assuming $\alpha < \beta$) if we take a value such that $\alpha < x < \beta$, then $x-\beta < 0$, therefore implying that $\frac{(x-\alpha)^3(x-\beta)}{x}<0$ or otherwise resulting $f(x)<n$, which is contradiction as we assumed $n$ is the least value of the function. Similarly, logic proceeds if the third power is on the other binomial.

 





\textbf{Solution 2 (AM-GM)}

Consider the function\begin{align*}f(x) = \frac{(x - 18)(x - 98)(x - 72)(x - k)}{x^2}.\end{align*}
Expanding this, we obtain\begin{align*}f(x) = \left( x + \frac{18 \times 98}{x} - (18 + 98)\right)\left( x + \frac{72k}{x} - (72 + k) \right).\end{align*}


 Let \( y = x + \frac{m}{x} \) (where \( x > 0 \)). By the AM-GM inequality \( a + b \geq 2\sqrt{ab} \), we have\begin{align*}x + \frac{m}{x} \geq 2\sqrt{m}.\end{align*}
Assuming \( x = \sqrt{m} \), the minimum value is \( 2\sqrt{m} \).

 Let \( y_1 = x + \frac{18 \times 98}{x}  - (18 + 98) \). Then,

 when\(x  = \sqrt{18 \times 98} = 42.\) We obtain \(y_{min} = 42+42-116 =-32\)

 Let \( y_2 = x + \frac{72 \times k}{x} - (72 + k) \). 

 When\(x  = \sqrt{72 \times k}\),We obtain \(y_{min} = 2(\sqrt{72 \times k})-(72+k)\).

 Since \( y_1 = y_2 \), we have\begin{align*}2(\sqrt{72 \times k})-(72+k) = -32,\end{align*}
which yields \( k = 8 \) or \( k = 200 \).

 With the same method, consider the function\begin{align*}f(x) = \underbrace{\left( x + \frac{18 \times 72}{x} - 90 \right)}_{y_1} \cdot \underbrace{\left( x + \frac{98k}{x} - (98 + k) \right)}_{y_2}.\end{align*}
When \( x_1 = \sqrt{18 \times 72} = 36 \),\begin{align*}y_1 = 2\times 36 - 90 = -18,\end{align*}
When \( x_2 = \sqrt{98 \times k}\),\begin{align*}y_2 = 2(\sqrt{98 \times k})-(98+k),\end{align*}
Thus, \(y_1 = y_2 \) gives \( 2(\sqrt{98 \times k})-(98+k) = -18 \), and the minimum value corresponds to \( k = 32 \) or \( k = 200 \).

 In summary, \( k_1 = 8 \), \( k_2 = 32 \), and \( k_3 = 200 \), with their sum being$8+32+200=\boxed{240}$.

 Comment: The function $f$ in the solution is different (the denominator $x$ is replaced with $x^2$). So this is another problem.Another issue is that, even if both $y_1$ and $y_2$ assume their minimal value (-18), their product is positive, not necessarily a minimum. Also, there is a problem that the person assumed multiple values of $x$ at the multiplicative of the function and yet assumed it to be a single variable function.

 -J.Z.

 





\textbf{Solution 3 (elegant polynomial)}

We first do the same thing as Solution 1 did, but in mathematical language.

 Lemma: for Polynomial $P(x)$, if $P(x_0) = 0$ and $P'(x_0) = 0$ both hold, then $x_0$ is a multiple root of $P(x)$.

 This lemma is obvious since every root of $P(x)$ occurs $\deg P - 1$ times in the derivative polynomial and any single root of $P(x)$ can't be a root of $P'(x)$.

 Simply name m as the minimum value of $f(x)$.

 Then $f(x) \geq m$ and the equal sign holds if and only of $x=x_1$ or $x=x_2$.

 Define $P(x)=(x-18)(x-72)(x-98)(x-k)-mx$. Obviously, $x_1$ and $x_2$ satisfy the two properties in the lemma. So we've got:
 \[(x-18)(x-72)(x-98)(x-k)-mx = (x-x_1)^2(x-x_2)^2\]Take $x=18, 72, 98, k$, we get:
 \[(18-x_1)(18-x_2)= \pm\sqrt{-m}\sqrt{18}\]and so on.So we can construct $Q(t)=(t^2-x_1)(t^2-x_2)-\sqrt{-m}t$.As what we have seen, $\pm\sqrt{18}, \pm\sqrt{72}, \pm\sqrt{98}, \pm\sqrt{k}$ are four roots of $Q(t) = 0$.So we've got:
 \[Q(t)=(t^2-x_1)(t^2-x_2)-\sqrt{-m}t=(t\pm\sqrt{18})(t\pm\sqrt{72})(t\pm\sqrt{98})(t\pm\sqrt{k})\]Comparing the cubic coefficient:
 \[\pm\sqrt{18}\pm\sqrt{72}\pm\sqrt{98}\pm\sqrt{k}=0\]So
 \[k=(\pm\sqrt{18}\pm\sqrt{72}\pm\sqrt{98})^2\]Now $8, 32, 200, 512$ are accesible in this form. The final task is to eliminate $512$. In this occasion the four roots are either $\sqrt{18}, \sqrt{72}, \sqrt{98}, -\sqrt{512}$ or $-\sqrt{18}, -\sqrt{72}, -\sqrt{98}, \sqrt{512}$. In either way the constant coefficient of $Q(t)$ is
 \[x_1x_2 = -\sqrt{18}\sqrt{72}\sqrt{98}\sqrt{512}\]But, $x_1$ and $x_2$ are both positive, so we end up with contradiction.

 $8+32+200=\boxed{240}$

 





\textbf{Solution 4}

What this problem is asking is to choose $k$ for which there exists a real $c$ with the property that 
 \[\frac{(x-18)(x-72)(x-98)(x-k)}{x}=c\] at exactly two positive reals $x$. Specifically, $c$ is the minimum value of the expression for positive reals $x$. This rewrites into $(x-18)(x-72)(x-98)(x-k)-cx=0$ having two roots. Because it is a quartic, then it is either the square of a quadratic or the product of a cubic and a linear. However, for the latter case, note that this essentially contradicts our minimality assumption. Therefore, it needs to be the square of a quadratic. We can write that squared quadratic as $x^2-mx+n$. We would then have the equation 
 \[(x-18)(x-72)(x-98)(x-k)-cx=\left (x^2-mx+n\right )^2.\]

 Firstly, compare the coefficients of $x^3$. In the left equation, we can see by vietas that the term is $-(k+188)x^3$, while in the right equation this is $-2mx^3$. We then obtain the relationship $2mx^3=(k+188)x^3\implies k=2m-188$. So it is equivalent to finding the values of $m$ that work. Next, we can start plugging in convenient values for $x$. For $x=18$, we find that $-18c=(324-18m+n)^2$. If we plug in $x=72$, we also get $-72c=(5184-72m+n)^2$. Similarly, $x=98$ gets us $-98c=(9604-98m+n)^2$. Notice how we have three variables in a three-way systems of equations, so it can be solved. Dividing the second equation by the first equation, we get a new equation 
 \[\frac{-72c}{-18c}=\frac{(5184-72m+n)^2}{(324-18m+n)^2},\] which simplifies to 
 \[4=\left (\frac{5184-72m+n}{324-18m+n}\right )^2.\] Taking the square root of both sides and doing a little rearranging, it simplifies into $5184-72m+n=\pm 2(324-18m+n)$. Proceeding similarly, we can divide the third equation by the first equation to get 
 \[\frac{-98c}{-18c}=\frac{(9604-98m+n)^2}{(324-18m+n)^2}\implies \frac{49}{9}=\left (\frac{9604-98m+n}{324-18m+n}\right )^2\implies 3(9604-98m+n)=\pm 7(324-18m+n).\]

 At this point, we have now simplified it slightly down into two variables and two equations. We need both equations to be true, although the plus and the minus cases for each equation aren't necessarily both true at the same time. For convenience, let us call the four equations 
 \[5184-72m+n=2(324-18m+n),\] 
 \[5184-72m+n=-2(324-18m+n),\] 
 \[3(9604-98m+n)=7(324-18m+n), \text{ and}\] 
 \[3(9604-98m+n)=-7(324-18m+n)\] as equations 1, 2, 3, and 4, respectively. 

 To make things simpler, we can simplify each equation down.

 In equation 1, we have $5184-72m+n=648-36m+2n$, so $36m+n=4536$.

 In equation 2, we have $5184-72m+n=-648+36m-2n$, so $108m-3n=5832\implies 36m-n=1944$.

 In equation 3, $28812-294m+3n=2268-126m+7n$, so $168m+4n=26544\implies 42m+n=6636$.

 Then in equation 4, $28812-294m+3n=-2268+126m-7n$, so $420m-10n=31080\implies 42m-n=3108$.

 Now, recall our earlier condition how both of those equations with the $\pm$ must be true. Then we have four cases of which are true: Equations 1 and 3, 1 and 4, 2 and 3, and 2 and 4.

 Case 1: Equations 1 and 3 are true

 We then have $36m+n=4536$ and $42m+n=6636$. Subtracting the equations gives $6m=2100\implies m=350$. However, it is obvious that $n$ is negative. But then this would force that one of the roots in the quadratic $x^2-mx+n$ are negative, which contradicts our initial assumption that it occurs in the positive reals $x$. This case doesn't provide any solutions.

 Case 2: Equations 1 and 4 are true

 It requires $36m+n=4536$ and $42m-n=3108$. Adding the two equations gives $78m=7644\implies m=98$. Then $n$ is positive, so this case does indeed work. Applying the relationship $k=2m-188$ from earlier gives $k=8$ as a possibility.

 Case 3: Equations 2 and 3 are true

 Then $36m-n=1944$ and $42m+n=6636$. Adding the two equations, $78m=8580\implies m=110$. $n$ is indeed positive, so this case provides $k=2\cdot 110-188=32$.

 Case 4: Equations 2 and 4 are true

 The two equations are $36m-n=1944$ and $42m-n=3108$. Subtracting the first from the second, we see that $6m=1164\implies m=194$. We find that $n$ is also positive here, so it works. Then it can be easily found that $k=2\cdot 194-188=200$.

 Therefore, our possible values of $k$ are $k=8,32,200$, so the answer is $\boxed{240}$.

 


 \href{https://mathproblemsolvingskills.wordpress.com/}{MathProblemSolvingSkills.com}

 


 




\end{enumerate}